

% \section{Relationship of Assumption \ref{asm:ZO} and \ref{asm:FO}}

% Below, we prove that under regular conditions on $f(\,\cdot\,)$ and $F(\,\cdot\,;\xi)$, we can get Assumption \ref{asm:FO} from Assumption \ref{asm:ZO}. 

% \begin{lem} \label{lem:cor-main-assmp}
	% Under Assumption \ref{asm:ZO}, if both $f(\,\cdot\,)$ and $F(\,\cdot\,;\xi)$ are differentiable everywhere, we have
	% \begin{enumerate}
		%     \item[(a)] $f(\,\cdot\,)$ is $L$-Lipschitz continuous. 
		%     \item[(b)]  $ \E_{\xi \sim \Xi}[\nabla F(\vx;\xi)] = \nabla f(\vx)$ and $\E_{\xi \sim \Xi} \Vert \nabla F(\vx;\xi) \Vert^2 \le L^2$.
		% % Assumption \ref{asm:ZO} implies Assumption \ref{asm:FO}  with $\vg(\vx;\xi) = \nabla F(\vx; \xi)$ and  $G=L$.
		% % % Under Assumption \ref{asm:main}, we have
		% \end{enumerate}
	% \end{lem}

% \begin{proof}
	% Below, we derive (a) and (b) from Assumption \ref{asm:ZO} separately. 
	% \begin{enumerate}
		%     \item[(a)] By Assumption \ref{asm:ZO}(a) and Jensen's inequality, for any $\vx, \vx' \in \sR^d$ we have that 
		%     \begin{align*}
			%         \vert f(\vx) - f(\vx') \vert \le \E_{\xi \sim \Xi} \vert f(\vx) - f(\vx') \vert \le \E_{\xi \sim \Xi}[L(\xi)] \Vert \vx - \vx' \Vert \le L \Vert \vx - \vx' \Vert. 
			%     \end{align*}
		%     \item[(b)] At any point $\vx \in \sR^d$ such that $\nabla f(\vx)$ and $\nabla F(\vx ;\xi)$ both exist, we have that
		%     \begin{align} \label{eq:lim-grad}
			%         \lim_{\vh \rightarrow \vzero} \frac{f(\vx+\vh) - f(\vx) - \E_{\xi \sim \Xi }\langle \nabla F(\vx;\xi), \vh \rangle  }{\Vert \vh \Vert} &= \lim_{\vh \rightarrow \vzero} \E_{\xi \sim \Xi } \left[ \frac{F(\vx+\vh; \xi) - F(\vx; \xi) - \langle \nabla F(\vx;\xi), \vh \rangle  }{\Vert \vh \Vert} \right].
			%     \end{align}
		%     The differentiability of $\nabla F(\,\cdot\,;\xi)$ at point $\vx $ implies that 
		%     \begin{align*}
			%         \lim_{\vh \rightarrow \vzero} \frac{F(\vx+\vh; \xi) - F(\vx; \xi) - \langle \nabla F(\vx;\xi), \vh \rangle  }{\Vert \vh \Vert}  = \vzero,
			%     \end{align*}
		%     while Assumption \ref{asm:ZO}(b) implies that 
		%     \begin{align*}
			%         \frac{\vert F(\vx+\vh; \xi) - F(\vx; \xi)\vert }{\Vert \vh \Vert} \le L(\xi), \quad {\rm and} \quad \frac{\vert \langle \nabla F(\vx;\xi), \vh \rangle \vert }{\Vert \vh \Vert}  \le L(\xi). 
			%     \end{align*}
		%     Since $\E[L(\xi)] \le L $, we can apply Lebesgue's dominated convergence theorem on (\ref{eq:lim-grad}) to obtain that 
		%     \begin{align*}
			%          \lim_{\vh \rightarrow \vzero} \frac{f(\vx+\vh) - f(\vx) - \E_{\xi \sim \Xi }\langle \nabla F(\vx;\xi), \vh \rangle  }{\Vert \vh \Vert}  = \vzero.
			%     \end{align*}
		%     This completes the proof of $ \E_{\xi \sim \Xi}[\nabla F(\vx;\xi)] = \nabla f(\vx)$. Finally, we calculate the second-order moment bound by
		%     \begin{align*}
			%         \E_{\xi \sim \Xi} \Vert \nabla F(\vx;\xi) \Vert^2 \le   \E_{\xi \sim \Xi} [L(\xi)]^2 \le L^2.
			%     \end{align*}
		% \end{enumerate}
	% \end{proof}

% \begin{rmk}
	% The condition that both $f(\,\cdot\,)$ and $F(\,\cdot\,;\xi)$ are differentiable everywhere in Lemma \ref{lem:cor-main-assmp} is mild and without loss of generality because the Rademacher’s theorem indicates that any Lipschitz continuous function is differentiable almost everywhere, and therefore such conditions can be fulfilled with probability 1 by adding arbitrarily small random permutations on the original functions. It is a standard technique used in previous works, which can refer to \textit{e.g.}  \citep[Proposition 2]{cutkosky2023optimal}, \citep[Lemma 2.3]{davis2022gradient}, or \citep[Lemma 3.2]{tian2022finite}.
	% \end{rmk}

\section{Proof of Randomized Smoothing}

\propfdeltasmooth*

\begin{proof}
	We verify each items one by one as follows.
	\begin{enumerate}
		\item[(a)] It follows the definition of $f_{\delta}$ (Definition \ref{dfn:f-delta}) and Jensen's inequality that 
		\begin{align*}
			\vert f_{\delta}(\vx) - f(\vx) \vert \le \E_{\vu \sim {\rm Unif}(\sB^d)} \vert f(\vx + \delta \vu) -f(\vx) \vert \le L \delta \E_{\vu \sim {\rm Unif}(\sB^d)} \Vert \vu \Vert \le L \delta.
		\end{align*}
		\item[(b)] By the definition of $f_{\delta}$ (Definition \ref{dfn:f-delta}), the Lipschitz continuity of $f$ and Jensen's inequality, for any $\vx, \vx' \in \sR^d$ we have that 
		\begin{align*}
			\vert f_{\delta}(\vx) - f_{\delta}(\vx') \vert &= \vert  \E_{\vu \sim {\rm Unif}(\sB^d)} [f(\vx+ \delta \vu)] -  \E_{\vu' \sim {\rm Unif}(\sB^d)} [f(\vx'+ \delta \vu')] \vert \\
			&= \vert  \E_{\vu \sim {\rm Unif}(\sB^d)} [f(\vx+ \delta \vu)] -  \E_{\vu \sim {\rm Unif}(\sB^d)} [f(\vx' + \delta \vu)] \vert \\
			&\le \E_{\vu \sim {\rm Unif}(\sB^d)} \vert f(\vx+ \delta \vu) - f(\vx' + \delta \vu) \vert \le L \Vert \vx - \vx' \Vert.
		\end{align*}
		\item[(c)] See \citep[Lemma 4]{kornowski2024algorithm} for a proof. 
		\item[(d)] It is first proven in \cite{yousefian2012stochastic}. See also \citep[Proposition 2.2]{lin2022gradient} for a more complete proof. 
	\end{enumerate}
\end{proof}

\section{Proof of the Nonconvex Optimistic Gradient Method}

% \lemsmoothonc*

% \begin{proof}
	% We start from the fundamental theorem of calculus that
	% \begin{align*}
		%     f(\vx_{t}) - f(\vx_{t-1}) =& \left \langle \int_0^1 \nabla f(\vx_{t-1}) + s \vDelta_t) {\rm d} s, \vDelta_t \right \rangle \\
		%     =& \langle \nabla f(\vy_t), \vDelta_t \rangle +  \left \langle \int_0^1 (\nabla f(\vx_{t-1}) + s \vDelta_t) - \nabla f(\vy_t) ) {\rm d} s, \vDelta_t \right \rangle
		% \end{align*}
	% For the right-hand side, we further apply
	% (1) Cauchy-Schwarz inequality, (2) $H$-smoothness of the function~$f$, (3) the identity $\vy_t = \vx_{t-1} + \vDelta_t/2$, and (4) $\Vert \Delta_t \Vert \le D$, which gives
	% \begin{align*}
		%     f(\vx_{t}) - f(\vx_{t-1}) \le& \langle \nabla f(\vy_t), \vDelta_t \rangle +  \int_0^1 \Vert\nabla f(\vx_{t-1}) + s \vDelta_t) - \nabla f(\vy_t) \Vert \Vert \vDelta_t \Vert {\rm d} s \\
		%     \le& \langle \nabla f(\vy_t), \vDelta_t \rangle +  H \int_0^1 \Vert\vx_{t-1} + s \vDelta_t-\vy_t \Vert \Vert \Vert \vDelta_t \Vert {\rm d} s  \\
		%     =& \langle \nabla f(\vy_t), \vDelta_t \rangle + H\int_0^1 \left \vert s- \frac{1}{2} \right \vert \Vert \vDelta_t \Vert^2 {\rm d} s = \langle \nabla f(\vy_t), \vDelta_t \rangle + \frac{H D^2}{4}. 
		% \end{align*}
	% Note that conditioned on $\Delta_t$, we have $\E [\mO_f(\vy_t)] =\nabla f(\vy_t)$. Therefore, we can take the expectation on both sides and then replace $\nabla f(\vy_t)$ with $\mO_f(\vy_t)$ on the right-hand side, which gives
	% \begin{align*}
		%     \E[f(\vx_{t}) - f(\vx_{t-1})] 
		%     \le& \E \langle \mO_f(\vy_t), \vDelta_t \rangle + \frac{H D^2}{4}. 
		% \end{align*}
	% Telescoping this inequality for $t = (k-1)M+1,\cdots,kM$ and rearranging the right-hand side, we obtain
	% \begin{align*}
		%     \E[f(\vx_{kM}) - f(\vx_{(k-1)M})] \le& \sum_{t=(k-1)M+1}^{k(k+1)M}  \E \langle \mO_f(\vy_t), \Delta_t - \vu_k \rangle + \sum_{t=(k-1)M+1}^{k(k+1)M}  \E \langle \mO_f(\vy_t), \vu_k \rangle + \frac{M H D^2}{4}.
		% \end{align*}
	% Letting $\vu_k = - D \sum_{t=(k-1)M+1}^{(k+1)M} \mO_f(\vy_t) / \Vert \sum_{t=(k-1)M+1}^{(k+1)M} \mO_f(\vy_t) \Vert $ in the above, we have
	% \begin{align*}
		%     \E[f(\vx_{kM}) - f(\vx_{(k-1)M})] \le& \sum_{t=(k-1)M+1}^{(k+1)M} \E \langle \mO_f(\vy_t), \Delta_t - \vu_k \rangle  - D \E \left \Vert \sum_{t=(k-1)M+1}^{(k+1)M} \mO_f(\vy_t) \right \Vert + \frac{M HD^2}{4}.
		% \end{align*}
	% Next, using the definition of $\sigma$-SFO that $\E \mO_f(\vy_t) = \nabla f(\vy_t)$ and $\E \Vert \mO_f(\vy_t) -\nabla f(\vy_t) \Vert^2 \le \sigma^2$, we know that $\E \left \Vert \sum_{t=(k-1)M+1}^{k(k+1)M} \mO_f(\vy_t) - \nabla f(\vy_t) \right \Vert \le \sqrt{M} \sigma$. Hence, using the triangle inequality, we can further obtain that
	% \begin{align*}
		%      \E[f(\vx_{kM}) - f(\vx_{(k-1)M})] \le& \sum_{t=(k-1)M+1}^{(k+1)M}  \E \langle \mO_f(\vy_t), \Delta_t - \vu_k \rangle  - D \E \left \Vert \sum_{t=(k-1)M+1}^{(k+1)M} \nabla f(\vy_t) \right \Vert + D\sqrt{M} \sigma+ \frac{MHD^2}{4}.
		% \end{align*}
	% Finally, summing up this inequality for $k = 1,\cdots,K$, dividing both sides by $KDM$, and rearranging gives
	% \begin{align*}
		%    \frac{1}{K} \E \left \Vert \frac{1}{M} \sum_{t=(k-1)M+1}^{(k+1)M} \nabla f(\vy_t) \right \Vert \le \frac{f(\vx_0) - f(\vx_{K})}{KDM} + \frac{{\rm Reg}_{T,K}}{KDM} + \frac{\sigma}{\sqrt{M}} + \frac{HD}{4}.
		% \end{align*}
	% \end{proof}

% \lemsmoothreg*

% \begin{proof}
	% Recall Lemma \ref{lem:odog} tells that
	% \begin{align} \label{eq:regret-odog-lem}
		% \begin{split}
			%     \sum_{t= (k-1)M +1}^{k M} \E \langle \mO_f(\vy_t), \vDelta_t - \vu_k \rangle \le& \frac{4 D^2}{\eta} + \frac{3\eta}{2}\E \Vert \mO_f(\vy_{(k-1)M}) - \mO_f(\vz_{(k-1)M}) \Vert^2 \\
			%     + & \sum_{t= (k-1)M +1}^{k M} \E \left[\eta \Vert \mO_f(\vy_t) - \mO_f(\vz_t) \Vert^2 - \frac{1}{4 \eta} \Vert \vDelta_t - \vDelta_{t-1} \Vert^2\right]. 
			% \end{split}
		% \end{align}
	% We know from  (1) Young's inequality, (2) the definition of $\sigma$-SFO, (3) and $H$-smoothness of $f$, and (4) the identity $\vy_t - \vz_t = (\vDelta_t - \vDelta_{t-1})/2$ that
	% \begin{align*}
		%     \Vert \mO_f(\vy_t) - \mO_f(\vz_t) \Vert^2 \le&
		% 3 \Vert \mO_f(\vy_t) - \nabla f(\vy_t) \Vert^2 + 3 \Vert \mO_f(\vz_t) - \nabla f(\vz_t) \Vert^2 + 3 \Vert \nabla f(\vy_t) - \nabla f(\vz_t) \Vert^3 \\
		% \le& 6 \sigma^2 + 3H^2 \Vert \vy_t - \vz_t \Vert^2 = 6 \sigma^2 + \frac{3H^2}{4} \Vert \vDelta_t - \vDelta_{t-1} \Vert^2,
		% \end{align*}
	% which is an analog to \citep[Lemma B.3]{patitucci2025improving}. And the remaining proofs are also similar to those in \citep{patitucci2025improving}. Substituting into the regret bound in inequality (\ref{eq:regret-odog-lem}), rearranging, and using the fact $\Vert \vDelta_t \Vert \le D$, we have 
	% \begin{align*}
		%   \sum_{t= (k-1)M +1}^{k M} \E \langle \mO_f(\vy_t), \vDelta_t - \vu_k \rangle \le& \frac{4 D^2}{\eta} + \frac{9\eta H^2 D^2}{2} + 9 M \eta \sigma^2 - \left( \frac{1}{4 \eta} - \frac{3 H^2}{4} \right)  \sum_{t= (k-1)M +1}^{k M} \E \Vert \vDelta_t - \vDelta_{t-1} \Vert^2.   
		%   % \\
		%   %   &\E [{\rm Reg}_{T,K}] \le \frac{4 KD^2}{\eta} +
		%   %   15 \eta KM \sigma^2 + \frac{15 \eta H^2 D^2}{8}  -
		%   %   \left( \frac{1}{4 \eta} -\frac{15 \eta H^2 }{8}   \right)\sum_{t=2}^{KM} \E  
		% \end{align*}
	% Note that by setting $\eta \le \frac{1}{\sqrt{3} H}$,  we can guarantee that $\frac{1}{4 \eta} -\frac{3 \eta H^2}{4} \ge 0$ and drop the last term in the above. Then, we obtain 
	% \begin{align*}
		%     \sum_{t= (k-1)M +1}^{k M} \E \langle \mO_f(\vy_t), \vDelta_t - \vu_k \rangle \le& \frac{4 D^2}{\eta} +\frac{9\eta H^2 D^2}{2} + 9 M \eta \sigma^2.   
		% \end{align*}
	% Next, by selecting the stepsize as $\eta = \min \left\{ \frac{1}{2H}, \frac{2D}{3 \sigma \sqrt{M}} \right\}$ and using the basic inequality $\frac{1}{\min\{a,b \}} \le a+b$, we can further bounded the right-hand side above as
	% \begin{align*}
		%  \sum_{t= (k-1)M +1}^{k M} \E \langle \mO_f(\vy_t), \vDelta_t - \vu_k \rangle \le 18 H D^2 + 12 \sigma \sqrt{M} D.
		% \end{align*}
	% Summing up this inequality for all $k = 1,\cdots,K$, we know that the $K$-shifting regret can be bounded by $K$ times of its right-hand side, which gives
	% \begin{align*}
		%     \E [{\rm Reg}_{T,K} ] \le 6K \left( 3 H D^2 +2 \sigma \sqrt{M} D^2 \right).
		% \end{align*}
	% \end{proof}
% \thmodogs*

% \begin{proof}
	% Plug the regret bound of Lemma \ref{lem:reg-smooth-surrogate} into Lemma \ref{lem:smooth-o2nc}, we get
	% \begin{align*}
		%     \E \Vert \nabla f(\bar \vy_{\rm out}) \Vert \le 
		%     \frac{F_0}{KDM} +
		%     \frac{18 H D}{M} + \frac{13 \sigma}{\sqrt{M}} + \frac{HD}{4}.
		% \end{align*}
	% For $D = \delta/M$, $T = KM$ and $H= c \sqrt{d} L / \delta$, we have that
	% \begin{align*}
		%      \E \Vert \nabla f(\bar \vy_{\rm out}) \Vert \le 
		%     \frac{M F_0}{\delta T} +
		%     \frac{73c \sqrt{d} L}{4 M} + \frac{13 \sigma}{\sqrt{M}}.
		% \end{align*}
	% By selecting $M = \max \left\{ \left( \frac{73 c \sqrt{d} L \delta T}{4 F_0} \right)^{1/2}, \left( \frac{13 \sigma \delta T}{F_0} \right)^{2/3} \right\}$ and using the basic inequality $\max\{a,b \} \le a+b$, we can further bounded the right-hand side above as
	% \begin{align*}
		%      \E \Vert \nabla f(\bar \vy_{\rm out}) \Vert \le \left( 
		%      \frac{73 c\sqrt{d} L F_0}{\delta T}
		%      \right)^{1/2} + \frac{2 F_0^{1/3}\left(13 \sigma \right)^{2/3}}{\left(\delta T \right)^{1/3}}.
		% \end{align*}
	% Therefore, to find a $(\delta,\epsilon)$-Goldstein stationary point, the number of oracles $\sO_{f_\delta}$ required by the algorithm is no more than
	% \begin{align*}
		%     T = \gO \left( \frac{\sqrt{d} L F_0}{\delta \epsilon^2} + \frac{\sigma^2 F_0}{\delta \epsilon^{3}} \right).
		% \end{align*}
	% \end{proof}
\lemregretong*

\begin{proof}
	We start from Lemma \ref{lem:odog} that
	\begin{align*}
		\E [{\rm Reg}_{T,K}] \le \frac{4 K D^2}{\eta} + \frac{5 \eta}{2} \sum_{t=1}^{KM} \E \Vert \mO(\vy_t) - \mO(\vy_{t-1}) \Vert^2.
	\end{align*}
	Applying Young's inequality on the second term, we further obtain that
	\begin{align*}
		\E [{\rm Reg}_{T,K}] \le& \frac{4 K D^2}{\eta} + \frac{15 \eta}{2} \sum_{t=1}^{KM} \E [\Vert \nabla f_\delta(\vy_t) - \nabla f_\delta(\vy_{t-1}) \Vert^2 \\
		+ &\Vert \mO(\vy_t) - \nabla f_\delta(\vy_t) \Vert^2 + \Vert \mO(\vy_{t-1}) - \nabla f_\delta(\vy_{t-1}) \Vert^2].
	\end{align*}
	By Proposition \ref{prop:f-delta-smooth}(d) we have $\Vert \nabla f_\delta(\vy_t) - \nabla f_\delta(\vy_{t-1}) \Vert \le (c \sqrt{d} L /\delta) \Vert \vy_t - \vy_{t-1} \Vert$, and by the definition of $(\delta,\sigma)$ oracle (Definition \ref{dfn:delta-SFO}) we have $\E \Vert \mO(\vy) - \nabla f_\delta(\vy) \Vert^2 \le \sigma^2 $. Therefore, we get
	\begin{align*}
		\E [{\rm Reg}_{T,K}] \le \frac{4 K D^2}{\eta} + \frac{15 \eta c^2 d L^2 }{2\delta^2}  \sum_{t=1}^{KM} \E \Vert \vy_t - \vy_{t-1} \Vert^2  + 15 \eta KM \sigma^2.
	\end{align*}
	In Algorithm \ref{alg:O2NC}, we have $\vy_t - \vy_{t-1} = (\vx_{t-1} + s_t \vDelta_t) - (\vx_{x-2} + s_{t-1} \vDelta_{t-1}) = (1- s_{t-1}) \vDelta_{t-1} + s_t \vDelta_t \in \sB_{3D}^d(\vzero) $. Therefore, we can further bound the right-hand side in the above inequality as
	\begin{align*}
		\E[{\rm Reg}_{T,K}] \le \frac{4 K D^2}{\eta} + 15 \eta KM \left( \frac{9 c^2 d L^2 D^2}{2 \delta^2} + \sigma^2 \right).
	\end{align*}
\end{proof}

\thmong*
\begin{proof}
	Selecting the stepsize $\eta = \Theta \left( \frac{\delta}{L\sqrt{d M}} \land \frac{D}{\sigma \sqrt{M}}  \right)$ to minimize the regret in Lemma \ref{lem:regret-o2ng}, we obtain
	\begin{align*}
		\E[{\rm Reg}_{T,K}] = \gO \left( K \left( \frac{D^2 L \sqrt{dM}}{\delta} + D \sigma \sqrt{M}  \right)  \right).
	\end{align*}
	Substituting the above regret bound into Lemma \ref{lem:O2NC}, we have
	\begin{align*}
		\E\left[{\rm dist} \left(\partial_\delta f_{\delta}(\bar \vy_{\rm out}), \vzero \right) \right] = \gO \left( \frac{f_\delta(\vzero) - f_\delta^*}{K DM} + \frac{\sqrt{d} LD}{\sqrt{M}} + \frac{\sigma}{\sqrt{M}} \right).
	\end{align*}
	Note that by Assumption \ref{asm:Lip} we have $f(\vzero) - f^* \le \gamma$ and by Proposition \ref{prop:f-delta-smooth}(a) we have $\vert f_\delta(\,\cdot\,) - f(\,\cdot\,) \vert \le \delta L$. Therefore, we have $f_\delta(\vzero) - f_\delta^* \le \gamma + \delta L$ and
	\begin{align*}
		\E\left[{\rm dist} \left(\partial_\delta f_{\delta}(\bar \vy_{\rm out}), \vzero \right) \right] = \gO \left( \frac{\gamma + \delta L}{K DM} + \frac{\sqrt{d} LD}{\delta \sqrt{M}} + \frac{\sigma}{\sqrt{M}} \right).
	\end{align*}
	For $T = KM$ and $D = \delta/ M$, we have
	\begin{align*}
		\E\left[{\rm dist} \left(\partial_\delta f_{\delta}(\bar \vy_{\rm out}), \vzero \right) \right] = \gO \left( \frac{M(\gamma + \delta L)}{\delta T} + \frac{\sqrt{d} L}{M^{3/2}} + \frac{\sigma}{\sqrt{M}} \right).
	\end{align*}
	Selecting the parameter $M = \Theta \left( \left( 
	\frac{\sqrt{d} L \delta T}{\gamma + \delta L}
	\right)^{2/5} \lor \left( \frac{\sigma \delta T}{\gamma + \delta L} \right)^{2/3} \right)$ to minimize the right-hand side, we obtain
	\begin{align*}
		\E\left[{\rm dist} \left(\partial_\delta f_{\delta}(\bar \vy_{\rm out}), \vzero \right) \right]  = \gO \left(
		d^{1/5} L^{2/5} \left( \frac{\gamma + \delta L}{\delta T} \right)^{3/5} + \sigma^{2/3} \left( \frac{\gamma + \delta L}{\delta T} \right)^{1/3}
		\right).
	\end{align*}
	Recalling Proposition that \ref{prop:f-delta-smooth}(d) $\partial_{\delta} f_{\delta}(\,\cdot\,) \in \partial_{2 \delta} f(\,\cdot\,)$, we know that the algorithm can provably find a $(2\delta,\epsilon)$-Goldstein stationary point in total iterations no more than 
	\begin{align*}
		T = \gO \left( \frac{\gamma + \delta L}{\delta}  \left( \frac{ d^{1/3} L^{2/3}}{\epsilon^{5/3}} + \frac{\sigma^2}{\epsilon^3} \right) \right).
	\end{align*}
\end{proof}
% \lemregret*
% \begin{proof}
	% Note that Line \ref{line:O2NC-extra} and \ref{line:O2NC-OGD} in Algorithm \ref{alg:fast-O2NC} are exactly the update of optimistic online gradient descent \citep{hazan2010extracting,rakhlin2013online} or extragradient \cite{korpelevich1976extragradient} on the domain $\sB(\vzero, D)$. Following the standard analysis for this update, (see \textit{e.g.} \citep[Lemma 15]{chen2021impossible}, for any $\vu \in \sB(\vzero, D)$, we have 
	% \begin{align*}
		%     \langle \vg_t , \vDelta_{t} - \vu \rangle &\le  \langle \vg_{t} - \vh_{t}, \vDelta_{t}  - \vv_t \rangle + \frac{1}{2 \eta} \Vert \vu - \vv_t \Vert^2 \\
		%     &- \frac{1}{2 \eta} \Vert \vu - \vv_{t+1} \Vert^2 - \frac{1}{2 \eta } \Vert \vv_{t+1} - \vDelta_{t}  \Vert^2 - \frac{1}{2 \eta} \Vert \vDelta_{t}  - \vv_t \Vert^2.
		% \end{align*}
	% We can use Young's inequality to get
	% \begin{align*}
		%     \langle \vg_t , \vDelta_{t}  - \vu \rangle &\le  \frac{\eta}{2} \Vert \vg_t - \vh_{t} \Vert^2 + \frac{1}{2 \eta} \Vert \vu - \vv_t \Vert^2 - \frac{1}{2 \eta} \Vert \vu - \vv_{t+1} \Vert^2 
		% \end{align*}
	% {\color{red} Use ODOG analysis to provide a tighter regret bound.}
	
	
	% Telescoping yields that
	% \begin{align*}
		%     R_M(\vu_1,\cdots,\vu_K) &= \sum_{k=1}^K \sum_{m= (k-1)M +1}^{k M} \langle \vg_m, \vv_{m+1/2} - \vu_k \rangle \\
		%     &\le \frac{\eta}{2} \sum_{k=1}^K \sum_{m = (k-1)M +1}^{kM} \Vert \vg_m - \vg_{m-1} \Vert^2 + \frac{1}{2 \eta} \sum_{k=1}^K  \Vert \vu_k  - \vv_{(k-1)M+1)} \Vert^2.  
		% \end{align*}
	% Taking expectation, and leveraging the variance bound in the $(\delta,\sigma)$ SFO (Definition \ref{dfn:SFO}), we have that
	% \begin{align*}
		%     \E[ R_M(\vu_1,\cdots,\vu_K) ] &\le \eta \sum_{k=1}^K \sum_{m = (k-1)M +1}^{kM} \E \Vert \nabla f_{\delta}(\vy_{m}) - \nabla f_{\delta}(\vy_{m-1}) \Vert^2  \\
		%     &\quad  +  4 \eta  KM \sigma^2 +  \frac{1}{2 \eta} \sum_{k=1}^K  \Vert \vu_k  - \vv_{(k-1)M+1)} \Vert^2, 
		% \end{align*}
	% where we define $\vy_0 = \vzero$ in the above.
	% Now, we use Proposition \ref{prop:f-delta-smooth}(b) that $f_{\delta}(\,\cdot\,)$ is $L$-Lipschitz and  Proposition \ref{prop:f-delta-smooth}(d) that $f_{\delta}(\,\cdot\,)$ is $(c \sqrt{d}  / \delta)$-smooth for some numerical constant $c>0$ to get
	% \begin{align} \label{eq:tele-regret}
		% \begin{split}
			% \E[ R_M(\vu_1,\cdots,\vu_K) ] &\le \frac{ c^2 \eta  d L^2}{\delta^2}  \sum_{k=1}^K \sum_{m = (k-1)M +1}^{kM} \Vert \vy_m - \vy_{m-1} \Vert^2   \\
			%         &\quad +  4 \eta  KM \sigma^2 +  \frac{1}{2 \eta} \sum_{k=1}^K  \Vert \vu_k  - \vv_{(k-1)M+1)} \Vert^2.
			% \end{split}
		% \end{align}
	% By the update of the algorithm and triangle inequality, we can obtain that
	% \begin{align*}
		%     \Vert \vy_t - \vy_{t-1} \Vert \le \Vert \vy_t - \vx_{t-1} \Vert + \Vert \vx_t - \vx_{t-1} \Vert + \Vert \vy_{t-1} - \vx_{t-1} \Vert \le 3D.
		% \end{align*}
	% In addition, we also have
	% \begin{align*}
		%     \Vert \vu_k  - \vv_{(k-1)M+1)} \Vert \le \Vert \vu_k \Vert + \Vert \vv_{(k-1)M+1)} \Vert \le 2D.
		% \end{align*}
	% Using these two bounds in (\ref{eq:tele-regret}), we then get
	% \begin{align*}
		%     \E[ R_M(\vu_1,\cdots,\vu_K) ] &\le \frac{9 c^2 \eta d L^2 K M  D^2}{\delta^2}  +  4 \eta  KM \sigma^2 + \frac{2K D^2}{\eta}. 
		% \end{align*}
	% \end{proof}

% \thmupper*
% \begin{proof}
	
	
	% Setting 
	% \begin{align} \label{eq:setting-eta}
		%     \eta = \sqrt{\frac{2 K D^2}{ 9 c^2 d L^2 KM D^2/ \delta^2  + 4 K M \sigma^2 }}
		% \end{align}
	% yields that
	% \begin{align*}
		%     \E[ R_M(\vu_1,\cdots,\vu_K) ] &\le \sqrt{2 K } D \cdot \sqrt{9 c^2 d L^2 KM D^2/ \delta^2  + 4 K M \sigma^2 } \\
		%     &\le 3 c \sqrt{2d M} L D^2 K / \delta + 2\sqrt{2 M } DK \sigma.
		% \end{align*}
	% We now substitute the above regret bound into (\ref{eq:regre-O2NC}) to get
	% \begin{align*}
		%      &\quad \frac{1}{K} \sum_{k=1}^{K-1} \E \left \Vert \frac{1}{M} \sum_{m=1}^M \nabla f_{\delta}( \vy_{(k-1)M+m} ) \right \Vert \\
		%      &\le \frac{f_{\delta}(\vx_0) - f_{\delta}^*}{ D T} + \frac{3 c \sqrt{2d M} L D K}{\delta T}  + \frac{2\sqrt{2 M } K \sigma}{T} + \frac{\sigma}{\sqrt{M}}.
		% \end{align*}
	% Next, we substitute $D =\delta / M$ and $T = KM$ to get
	% \begin{align*}
		%     &\quad \frac{1}{K} \sum_{k=1}^{K-1} \E \left \Vert \frac{1}{M} \sum_{m=1}^M \nabla f_{\delta}( \vy_{(k-1)M+m} ) \right \Vert \\
		%     &\le \frac{M (f_{\delta}(\vx_0) - f_{\delta}^*)}{\delta T} + \frac{3 c \sqrt{2 d} L }{M^{\frac{3}{2}}} + \frac{(2 \sqrt{2} + 1) \sigma}{\sqrt{M}}.
		% \end{align*}
	% Furthermore, by Proposition \ref{prop:f-delta-smooth} we know $\vert f_{\delta}(\,\cdot\,) - f(\,\cdot\,) \vert \le \delta L$, and hence $f_{\delta}(\vx_0) - f_{\delta}^* \le \gamma + \delta L$ in the above.
	% From Proposition \ref{prop:Q-FO-estimator} we know a first-order quantum algorithm with $\tilde \gO(B)$ QSFO complexity can ensure $\sigma \le \sqrt{d} L / B)$, and from Proposition \ref{prop:Q-ZO-estimator} we know a zeroth-order quantum algorithm with $\tilde \gO(B)$ QSZO complexity can ensure $\sigma \le d L / B$. 
	% Here, we can view $B$ as a virtual batch size. To show the upper bound for both first-order and zeroth-order in a unified way, we define $\sigma = d^{\alpha} L / B$, where $\alpha = 1/2$ for first-order method and $\alpha = 1$ for zeroth-order method. 
	% We define
	% the total complexity of the algorithm (zeroth-order or first-order) by $N = T B$. 
	% Therefore, the above inequality can also be written as
	% \begin{align*}
		%     &\quad \frac{1}{K} \sum_{k=1}^{K-1} \E \left \Vert \frac{1}{M} \sum_{m=1}^M \nabla f_{\delta}( \vy_{(k-1)M+m} ) \right \Vert \\
		%     &\le \frac{B M (\gamma + \delta L)}{\delta N} + \frac{3 c \sqrt{2 d} L }{M^{\frac{3}{2}}} + \frac{(2 \sqrt{2} + 1) d^{\alpha} L}{\sqrt{M} B }.
		% \end{align*}
	% From the above inequality, we can see the optimal choice of $B$ is 
	% \begin{align} \label{eq:setting-B}
		%     B = \frac{ (2 \sqrt{2} + 1)^{\frac{1}{2}} d^{\frac{\alpha}{2}} L^{\frac{1}{2}} (\delta N)^{\frac{1}{2}} }{ M^{\frac{3}{4} } (\gamma + \delta L)^{\frac{1}{2}}}.
		% \end{align}
	% The above setting yields that
	% \begin{align*}
		%      &\quad \frac{1}{K} \sum_{k=1}^{K-1} \E \left \Vert \frac{1}{M} \sum_{m=1}^M \nabla f_{\delta}( \vy_{(k-1)M+m} ) \right \Vert \\ 
		%      &\le \frac{ (2 \sqrt{2} + 1)^{\frac{1}{2}} M^{\frac{1}{4}} (\gamma +\delta L)^{\frac{1}{2}} d^{\frac{\alpha}{2}} L^{\frac{1}{2}}}{ (\delta N)^{\frac{1}{2}} } + \frac{3 c \sqrt{2 d} L }{M^{\frac{3}{2}}}.
		% \end{align*}
	% From the above, we can see the optimal choice of parameter $M$ is
	% \begin{align}  \label{eq:setting-M}
		%     M = \left( \frac{3 c \sqrt{2 d} L (\delta N)^{\frac{1}{2}}}{ (2 \sqrt{2} + 1)^{\frac{1}{2}} M^{\frac{1}{4}} (\gamma +\delta L)^{\frac{1}{2}} d^{\frac{\alpha}{2}} L^{\frac{1}{2}} }  \right)^{\frac{4}{7}},
		% \end{align}
	% which will lead to the final upper bound of
	% \begin{align*}
		%      &\quad \frac{1}{K} \sum_{k=1}^{K-1} \E \left \Vert \frac{1}{M} \sum_{m=1}^M \nabla f_{\delta}( \vy_{(k-1)M+m} ) \right \Vert \\  
		%      &\le  \left( \frac{ 3c \sqrt{2} (2 \sqrt{2} +1)^3 (\gamma+ \delta L)^3 d^{3 \alpha+\frac{1}{2}} L^4 }{ (\delta N)^3 } \right)^{\frac{1}{7}}.
		% \end{align*}
	% It means the algorithm will provably output a $(\delta,\epsilon)$-stationary point in the total complexity of  
	% \begin{align*}
		%     N = \gO\left(  \frac{d^{\alpha+ \frac{1}{6}} L^{\frac{4}{3}}(\gamma + \delta L)}{\delta \epsilon^{\frac{7}{3}}}  \right).
		% \end{align*}
	% \end{proof}
% \paragraph{Upper Bound of (a)} From the previous discussion, we know that the output of $\mA_K$ is independent of $\vu_T$ and (\ref{eq:bound-uT}) holds conditioned on $\mA_K$ under event $G$. Therefore, for any $\vx \sim p_{K \mid G}$, for $\vz = \rho_R(\vx)$ with $R = 2 \sqrt{T}$ we have that $\vert \langle \vu_T, \vz \rangle \vert < 1 /4$. Our next goal is to following similar analysis to \citep[Lemma 25]{cutkosky2023optimal} to show that this implies that $\vx$ can not lie in the sets of $(1/2)$-stationary points $S =  \{ \vx \in \sR^d : \Vert \nabla \hat f_{T;\mU}(\vx) \Vert \le 1/2 \}$. 
% We consider two cases, either $\Vert \vx \Vert > \sqrt{T}$ or not. 

% If $\Vert \vx \Vert > \sqrt{T}$, we have that
% \begin{align*}
	%     \Vert \nabla \hat f_{T; \mU} (\vx) \Vert &\ge \lambda \Vert \nabla q_R(\vx) \Vert - \Vert \mJ_{\rho_R} (\vx) \Vert \Vert \nabla \bar f_{T; \mU}(\vy) \Vert \\
	%     &\ge \lambda \Vert \nabla q_R(\vx) \Vert - \Vert \nabla \bar f(\vy) \Vert \\
	%     &\ge \frac{\lambda \Vert \vx \Vert}{\sqrt{1 + \Vert \vx \Vert^2/ R^2}} - L_0 \sqrt{T} \\
	%     &\ge 3 \lambda R - L_0 \sqrt{T} \ge .
	% \end{align*}








% Above, the first inequality comes from the triangle inequality; the second one uses that fact $\Vert \mJ_{\rho_R}(\vx) \Vert \le 1$ by (\ref{eq:Jacobian}); the third one uses Lemma \ref{lem:prop-q}(b) and Lemma \ref{lem:Carmon}(c); and the last line uses 
% the monotonicity of function $q_R(\,\cdot\,)$.
% \begin{align*}
	%      \vert \le \Vert \vu_T \Vert \Vert \vz \Vert \le \beta / 4. 
	% \end{align*}


% To further upper bound (\ref{eq:P-TV}), we need to bound $\BP( \vx \in S \mid \vx \sim p_{K \mid F}, F  )$. 
% It means that conditioned on both events $E$ and $F$, for any $\vx \sim p_{K \mid F}$, we have that
% \begin{align*}
	%     \vert 
	% \end{align*}
% Conditioned on event $E$, we have that $\vert \vu_T, \vx $

% the worst-case construction of QSFO in  is independent of $\mU_T$. Therefore, by choosing $\vu_T$ 


% % represent the states $\psi_K$ and $\psi_0$ conditioned on event $F$.

% % Then, by Markov's inequality, there exists an event $F$ with $\BP(F) \ge 2/3$ such that
% % \begin{align} \label{eq:norm-Ak-A0}
	% %     \Vert  \psi_{K \mid F}  - \psi_{0 \mid F} \Vert^2 \le \frac{1}{48 T^2},
	% % \end{align}

% % where we use 
% % Next, conditioned on event $F$,  we let $p_{0 \mid F} = \vert \psi_K \vert^2$ and $p_{K \mid F} = \vert \psi_0 \vert^2$










% to denote the  $\psi_{0 \mid F}$ and $\psi_{K \mid F}$. 


% Let $F$ be the event that the above happens, then we have $\BP(F) \ge 2/3$.

% Conditioned on event $F$, 
% where  Define $\psi_{K \mid F}(\vx)$ and $\psi_{0 \mid F}(\vx)$ the wave functions of these states, which satisfy that $\vert \psi_{0 \mid F} (\vx) \vert^2 = p_{0 \mid F}$ and   $\vert \psi_{K \mid F} (\vx) \vert^2 = p_{K \mid F}$.


% We define $\psi_{K \mid \mU, \zeta }$ and $\psi_{0 \mid \mU, \zeta }$ to 


% Define
% $p_{K}(\vx) $ and $p_0(\vx)$ as the probability distribution obtained by measuring the state $\psi_K $ and $\psi_0 $, respectively. Similarly, we can define the conditional distribution $p_{K; \mU, \zeta}$ and $p_{0; \mU, \zeta}$, and we have the relationship that $p_{K;\mU,\zeta}(\vx) = \vert \psi_{K;\mU,\zeta}(\vx) \vert^2$ and $p_{0;\mU,\zeta}(\vx) = \vert \psi_{0;\mU,\zeta}(\vx) \vert^2$.
%  , we have that


% \begin{align*}
	%     d_{\rm TV}(p_0,p_K) &= \frac{1}{2}\iiint  \vert p_0(\vx) - p_K(\vx) \vert {\rm d} \vx {\rm d} \mU {\rm d} \zeta  \\
	%     &=  \frac{1}{2}  \iint_F \int \vert p_0(\vx) - p_K(\vx) \vert {\rm d} \vx {\rm d} \mU {\rm d} \zeta +  \frac{1}{2} \iint_{F^C} \int  \vert p_0(\vx) - p_K(\vx) \vert {\rm d} \vx {\rm d} \mU {\rm d} \zeta \\
	%     &= \frac{1}{2}  \iint_F \int \vert p_{0; \mU, \zeta}(\vx) p_\mU(\mU) p_\zeta(\zeta) - p_{K; \mU, \zeta}(\vx) p_\mU(\mU) p_\zeta(\zeta) \vert {\rm d} \vx {\rm d} \mU {\rm d} \zeta \\
	%     &\quad + \frac{1}{2}  \iint_{F^C} \int \vert p_{0; \mU, \zeta}(\vx) p_\mU(\mU) p_\zeta(\zeta) - p_{K; \mU, \zeta}(\vx) p_\mU(\mU) p_\zeta(\zeta) \vert {\rm d} \vx {\rm d} \mU {\rm d} \zeta \\
	%     &\le  \frac{1}{2}  \iint_F \int \vert p_{0; \mU, \zeta}(\vx) p_\mU(\mU) p_\zeta(\zeta) - p_{K; \mU, \zeta}(\vx) p_\mU(\mU) p_\zeta(\zeta) \vert {\rm d} \vx {\rm d} \mU {\rm d} \zeta + \sR() \\
	%     &= \frac{1}{2}\int \left\vert \vert \psi_0(\vx) \vert^2 - \vert \psi_K(\vx) \vert^2 \right \vert {\rm d} \vx \\
	%     &= \frac{1}{2}\int \left(\vert \psi_0(\vx) \vert + \vert \psi_K(\vx) \vert \right) \cdot \left \vert \psi_0(\vx) - \psi_K(\vx) \right \vert {\rm d} \vx \\
	%     &\le \frac{1}{2} \sqrt{ \int  \left(\vert \psi_0(\vx) \vert + \vert \psi_K(\vx) \vert \right)^2 {\rm d} \vx} \sqrt{ \int  \left \vert \psi_0(\vx) - \psi_K(\vx) \right \vert^2 {\rm d} \vx} \\
	%     &\le \frac{1}{2} \sqrt{ 2 ( \Vert \psi_0 \Vert^2 + \Vert \psi_K \Vert^2)  } \cdot  \Vert \psi_0 - \psi_K \Vert  = \Vert \psi_0 - \psi_K \Vert \le \frac{1}{6T},
	% \end{align*}

% %where the randomness comes from both $\mU$ and the stochasticity of QSFO. We let the above randomness be included by a random variable $\zeta$.

% The total variation distance between these two distributions is
% \begin{align*}
	%     d_{\rm TV}(p_K, p_0) = 
	% \end{align*}
% % and the quantum oracle $\hat \mO_{T;\mU}  \leftrightarrow \bar \mO_{T;\mU} $  that a quantum algorithm $\hat \mA_{\rm quan}$ which has access to oracle $\hat \mO_{T;\mU}$ can also viewed as an 

\section{Proof of Classical Gradient Estimators}
\propzoest*

% \subsection{Proof of Lemma \ref{lem:ZO-estimator}}
\begin{proof}
	To prove this proposition, it suffices to show that each single-batch estimator $\mO_j$ is a $(\delta, c \sqrt{d} L)$-SFO of the function $f$, which implies that the mini-batch estimator $\mO(\vx) = \frac{1}{B} \sum_{j=1}^B \mO_j(\vx)$ satisfies $\E [\mO(\vx)] = \nabla f_{\delta}(\vx) $ and ${\rm Var}[\mO(\vx)] = \frac{1}{B} {\rm Var}[\mO_1(\vx)] \le \sigma^2$. In the following, we prove the unbiasedness and variance bound in Definition~\ref{dfn:delta-SFO} separately. 
	% For the simplicity of notation.
	% In the following we use $\E[\,\cdot\,]$ to denote $\E_{\vv \sim {\rm Unif}(\sS^{d-1}), \xi \sim \Xi}[\,\cdot\,]$ 
	\begin{enumerate}
		\item \underline{Unbiasedness.} See \citep[Lemma 1]{flaxman2005online} as an application of generalized Stokes' theorem.
		\item \underline{Variance Bound.} It is first proven in \citep[Lemma 10]{shamir2017optimal}, but without explicit numerical constants $c$. 
		A value of the constant $c = 4 \sqrt[4]{2} \sqrt{\pi} $ is given by \citep[Lemma D.1]{lin2022gradient}.
	\end{enumerate}
	
	
	\underline{}{} 
\end{proof}

% \subsection{Proof of Lemma \ref{lem:FO-estimator}}
\propfoest*
\begin{proof}
	Similar to the proof of Proposition \ref{prop:ZO-estimator}, it suffices to show that each single-batch estimator $\mO_j$ is a $(\delta, c \sqrt{d} L)$-SFO of the function $f$. Below, we prove the unbiasedness and variance bound in Definition~\ref{dfn:delta-SFO} separately.
	We use $\E[\,\cdot\,]$ to denote $\E_{\vu \sim {\rm Unif}(\sB^d), \xi \sim \Xi}[\,\cdot\,]$ for simplicity of notation.
	\begin{enumerate}
		\item \underline{Unbiasedness.}
		Since $F(\,\cdot\,,\xi)$ is Lipschitz continuous, Rademacher's Theorem guarantees that it is differentiable almost everywhere with respect to the Lebesgue measure. This ensures that $\nabla F(\vx+ \delta \vu; \xi)$ exist for almost all $\vu \in {\rm Unif}(\sB^d)$.
		Then, under the condition that $\nabla F(\vx+\delta \vu,;\xi)$ exists, we can show that the expectation jointly in $\xi, \vu$ and derivative are interchangeable, which implies that $ \E [\nabla F(\vx+ \delta \vu;\xi)] = \nabla f_\delta(\vx)$.
		Since we have $f_\delta(\vx) = \E[F(\vx+\delta\vu;\xi)]$ by definition, we know that the following identity holds:
		\begin{align} \label{eq:lim-grad}
			\begin{split}
				&\lim_{\vh \rightarrow \vzero} \frac{f_\delta(\vx+\vh) - f_\delta(\vx) - \E\langle \nabla F(\vx+\delta \vu;\xi), \vh \rangle  }{\Vert \vh \Vert} \\
				=& \lim_{\vh \rightarrow \vzero} \E\left[ \frac{F(\vx+\delta\vu+ \vh; \xi) - F(\vx+\delta \vu; \xi) - \langle \nabla F(\vx+\delta \vu;\xi), \vh \rangle  }{\Vert \vh \Vert} \right].    
			\end{split}
			% \begin{split}
				%      & \lim_{\vh \rightarrow \vzero}  \left[ \frac{f(\vx+\vh) - f(\vx) - \E_{ \xi \sim \Xi} \langle \nabla F(\vx + \delta \vu;\xi), \vh \rangle  }{\Vert \vh \Vert} \right] \\
				%     = &\lim_{\vh \rightarrow \vzero} \E_{ \vu \sim {\rm Unif}(\sB^d)}  \left[ \frac{F(\vx+\delta \vu +\vh) - F(\vx+ \delta \vu) - \langle \nabla F(\vx + \delta \vu;\xi), \vh \rangle  }{\Vert \vh \Vert} \right].
				% \end{split}
		\end{align}
		%  Without loss of generality, we assume that both $f(\,\cdot\,)$ and $F(\,\cdot\,;\xi)$ are differentiable everywhere.
		% Otherwise, the Rademacher’s theorem indicates that any Lipschitz continuous function is differentiable almost everywhere and therefore such conditions can be fulfilled with probability 1 for all the points queried by the algorithms after adding arbitrarily small random permutations on the original functions. It is a standard technique used in previous works, \textit{e.g.}  \citep[Proposition 2]{cutkosky2023optimal}, \citep[Lemma 2.3]{davis2022gradient}, or \citep[Lemma 3.2]{tian2022finite}. 
		% At any point $\vx \in \sR^d$ such that $\nabla f(\vx)$ and $\nabla F(\vx ;\xi)$ both exist, we have that
		% \begin{align} \label{eq:lim-grad}
			%     \lim_{\vh \rightarrow \vzero} \frac{f(\vx+\vh) - f(\vx) - \E_{\xi \sim \Xi }\langle \nabla F(\vx;\xi), \vh \rangle  }{\Vert \vh \Vert} &= \lim_{\vh \rightarrow \vzero} \E_{\xi \sim \Xi } \left[ \frac{F(\vx+\vh; \xi) - F(\vx; \xi) - \langle \nabla F(\vx;\xi), \vh \rangle  }{\Vert \vh \Vert} \right].
			% \end{align}
		The differentiability of $\nabla F(\,\cdot\,;\xi)$ at the point $\vx+ \delta \vu $ implies that 
		\begin{align*}
			\lim_{\vh \rightarrow \vzero} \frac{F(\vx+\delta \vu + \vh; \xi) - F(\vx+\delta \vu; \xi) - \langle \nabla F(\vx+\delta \vu;\xi), \vh \rangle  }{\Vert \vh \Vert}  = 0,
		\end{align*}
		while Assumption \ref{asm:Lip} implies that 
		\begin{align*}
			\frac{\vert F(\vx+\delta \vu+\vh; \xi) - F(\vx+\delta \vu; \xi)\vert }{\Vert \vh \Vert} \le L(\xi), \quad {\rm and} \quad \frac{\vert \langle \nabla F(\vx+\delta \vu;\xi), \vh \rangle \vert }{\Vert \vh \Vert}  \le L(\xi). 
		\end{align*}
		Since $\E[L(\xi)] \le L $, we can apply Lebesgue's dominated convergence theorem on (\ref{eq:lim-grad}) to obtain that 
		\begin{align*}
			\lim_{\vh \rightarrow \vzero} \frac{f_\delta(\vx+\vh) - f_\delta(\vx) - \E\langle \nabla F(\vx+\delta \vu;\xi), \vh \rangle  }{\Vert \vh \Vert} =0.
		\end{align*}
		This completes the proof of $ \E [\nabla F(\vx+ \delta \vu;\xi)] = \nabla f_\delta(\vx)$.    
		%     Therefore, we have
		% \begin{align*}
			%     \E [\nabla F(\vx+\delta \vu;\xi)] =  \E_{\vu} [\nabla f(\vx+ \delta \vu)] = \nabla \E_{\vu} [f(\vx + \delta \vu) ]= \nabla f_{\delta}(\vx).
			% \end{align*}
		% By Proposition \ref{prop:f-delta-smooth}(d) we know that the smoothed surrogate is differentiable everywhere, therefore $\nabla f_{\delta}(\vx)$ always exists. Next, by the Rademacher’s theorem $F(\,\cdot\,; \xi)$ is differentiable almost everywhere, therefore $\nabla F(\vx+ \delta \vu;\xi)$ with $\vu \sim {\rm Unif}(\sR^d)$ exists with probability 1 for any random index $\xi$. 
		% The remaining steps are identical to the proof of Lemma \ref{lem:cor-main-assmp}(e), except that now we need to take expectation \textit{jointly} on random variables $\vu$ and $\xi$, while Lemma \ref{lem:cor-main-assmp}(e) only takes expectation on $\xi$. 
		\item \underline{Variance Bound.} We directly bound the variance with the second-order moment and obtain
		\begin{align*}
			{\rm Var}[ \nabla F(\vx+ \delta \vu; \xi)] \le 
			\E \Vert \nabla F(\vx + \delta \vu;\xi) \Vert^2 \le \E \Vert \nabla F(\vx+ \delta \vu ; \xi) \Vert^2 \le \E[L(\xi)^2] \le L^2.
		\end{align*}
	\end{enumerate}
	
\end{proof}

\propdist* 

\begin{proof}
	As the proofs of the single-agent estimator, we complete the proof by verifying the unbiasedness and variance bound in Definition~\ref{dfn:delta-SFO} separately. Below, we use $\E[\,\cdot\,]$ to denote $\E_{\vu \sim {\rm Unif}(\sB^d), \xi_1 \sim \Xi_1,\cdots,\xi_m \sim \Xi_m}[\,\cdot\,]$ for the simplicity of notation.
	\begin{enumerate}
		\item \underline{Unbiasedness.} It follows from the interchangeability of mean, derivative, and expectation:
		\begin{align*}
			\E[\mO(\vx)] =&  \frac{1}{m} \sum_{i=1}^m \E [ \mO_i(\vx)] =   \frac{1}{m} \sum_{i=1}^m \nabla \E[ f_i( \vx + \delta \vu)]  \\
			=&\E \left[  \frac{1}{m}  \sum_{i=1}^m  \nabla f_i(\vx+ \delta \vu) \right] = \nabla \E[f(\vx+\delta \vu)] = \nabla f_\delta(\vx).
		\end{align*}
		\item \underline{Variance bound.} Using the identity $\E \Vert \sum_{i=1}^m \va_i \Vert^2 =  \sum_{i=1}^m \E \Vert \va_i \Vert^2$, we have
		\begin{align*}
			{\rm Var}\left[ \frac{1}{m} \sum_{i=1}^m \mO_i(\vx) \right] = \frac{1}{m^2} \sum_{i=1}^m {\rm Var}[\mO_i(\vx)] \le \frac{\sigma^2}{m}.
		\end{align*}
	\end{enumerate}
\end{proof}

\section{Proof of Quantum Gradient Estimators}

% Our quantum estimators rely on the following facts in the classical setting. 

%\subsection{Proof of Proposition \ref{prop:f-delta-smooth}}




% We provide the full details here for the sake of completeness.  Similar to the proof of Lemma \ref{lem:cor-main-assmp}(e), we start from 
% \begin{align} \label{eq:lim-grad-new}
	% \begin{split}
		%      &\quad  \lim_{\vh \rightarrow \vzero} \frac{f_{\delta}(\vx+\vh) - f_{\delta}(\vx) - \E \langle \nabla F(\vx + \delta \vu;\xi), \vh \rangle  }{\Vert \vh \Vert} \\
		%        &= \lim_{\vh \rightarrow \vzero} \E \left[ \frac{F(\vx+ \delta \vu + \vh; \xi) - F(\vx+ \delta \vu; \xi) - \langle \nabla F(\vx + \delta \vu;\xi), \vh \rangle  }{\Vert \vh \Vert} \right].
		% \end{split}
	% \end{align}
%     The existence of $\nabla F(\vx+ \delta \vu,;\xi)$ implies that 
%     \begin{align*}
	%         \lim_{\vh \rightarrow \vzero} \frac{F(\vx+\delta \vu+ \vh; \xi) - F(\vx+ \delta \vu; \xi) - \langle \nabla F(\vx+ \delta \vu;\xi), \vh \rangle  }{\Vert \vh \Vert}  = \vzero,
	%     \end{align*}
%     while Assumption \ref{asm:main}(c) implies that 
%     \begin{align*}
	%         \frac{\vert F(\vx+\vh; \xi) - F(\vx; \xi)\vert }{\Vert \vh \Vert} \le L(\xi), \quad {\rm and} \quad \frac{\vert \langle \nabla F(\vx;\xi), \vh \rangle \vert }{\Vert \vh \Vert}  \le L(\xi). 
	%     \end{align*}
%     Since $\E[L(\xi)] \le L $, we can apply Lebesgue's dominated convergence theorem on (\ref{eq:lim-grad}) to obtain that 
%     \begin{align*}
	%          \lim_{\vh \rightarrow \vzero} \frac{f(\vx+\vh) - f(\vx) - \E_{\xi \sim \Xi }\langle \nabla F(\vx;\xi), \vh \rangle  }{\Vert \vh \Vert}  = \vzero.
	%     \end{align*}
%     This completes the proof.

% Finally, the above analysis 
% fulfills the conditions in Lemma \ref{lem:cor-main-assmp}(e), which indicates that 
% \begin{align*}
	%     \E_{\vu \sim {\rm Unif}(\sB^d), \xi \sim \Xi} [\nabla F(\vx+\delta \vu;\xi)] = \nabla 
	% \end{align*}






% Now, we prove Proposition \ref{prop:Q-ZO-estimator} and \ref{prop:Q-FO-estimator} in a unified manner. 

We use Proposition \ref{prop:QVR} to both prove Proposition \ref{prop:Q-ZO-estimator} and \ref{prop:Q-FO-estimator} in a similar manner. Moreover, in the following proofs, we also give explicit constructions of the quantum gradient estimators for our algorithms.

% \subsection{Proof of Proposition \ref{prop:Q-ZO-estimator}}

\propqzoest*
\begin{proof}
	To invoke Proposition \ref{prop:QVR}, we define the random variable $Z_x$ as
	\begin{align*}
		Z_x =  \frac{d}{2 \delta} ( F(\vx + \delta \vv;\xi ) - F(\vx - \delta \vv;\xi) \vv, \quad {\rm where} \quad \vv \sim {\rm Unif}(\sS^{d-1}), ~~ \xi \sim \Xi.
	\end{align*}
	And in the proof of Proposition \ref{prop:ZO-estimator}, we have shown that $ \E[Z_x] = \nabla f_{\delta}(\vx)$ and ${\rm Var}[Z_x] \le 16 \sqrt{2 \pi} dL^2$. 
	Then by Proposition \ref{prop:QVR} the \texttt{QuantumMeanEstimation+} can output a $(\delta,\sigma)$-SFO oracle with $\tilde \gO( d L / \sigma)$ quantum sampling oracle of the random variable $Z_x$. Finally, we finish the proof by explicitly constructing the quantum sampling oracle. It is clear that it can be implemented by QSZO oracle (Definition \ref{dfn:QSZO}) given we have access to the quantum sampling oracle to $\vv$ and $\xi$. 
	The quantum sampling oracle to $\xi \sim \Xi$ is nothing else but a uniform distribution $\int_0^1 | \xi \rangle {\rm d} \xi$, which can be prepared using Hadamard gates. To construct the quantum sampling oracle to $\vv \sim {\rm Unif}(\sB^{d-1})$, we can mimic the Box-Muller transform in quantum computers, such as  \citep[Section 5.3]{chakrabarti2023quantum}. Specifically, we can first prepare two uniformly distributed random variables $\int_0^1 | \xi_1 \rangle {\rm d} \xi_1$ and $\int_0^1 | \xi_2 \rangle {\rm d} \xi_2$. Then applying the quantum Box-Muller transform given by
	\begin{align*}
		| \xi_1 \rangle \otimes |\xi_2 \rangle  \mapsto \left| \sqrt{-4 \ln \xi_1} \cos 2 \pi \xi_2 \right \rangle \otimes \left| \sqrt{-4 \ln \xi_1} \sin 2 \pi \xi_2 \right \rangle
	\end{align*}
	yields the following result:
	\begin{align*}
		\int_0^1 | \xi_1 \rangle {\rm d} \xi_1 \otimes \int_0^1 | \xi_2 \rangle {\rm d} \xi_2 \mapsto \int_{\sR} \frac{1}{\sqrt[4]{2 \pi}} \exp(-z_1^2 /4) | z_1 \rangle {\rm d} z_1 \otimes \int_{\sR} \frac{1}{\sqrt[4]{2 \pi}} \exp(-z_2^2 /4) | z_2 \rangle {\rm d} z_2.
	\end{align*}
	By applying the above procedure on $d/2$ pairs, we can obtain the quantum sampling oracle for $d$-dimensional normal distribution
	\begin{align*}
		\int_{\sR^d} \frac{1}{\sqrt[4]{2 \pi}} \exp( - \Vert \vz \Vert^2/4) | \vz \rangle {\rm d} \vz. 
	\end{align*}
	Next, we can apply a normalization transform $ | \vz \rangle \mapsto | \vz / \Vert \vz \Vert \rangle $ to the above state, by the isotropy of $d$-dimensional normal distribution, we get the uniform distribution over the unit sphere
	\begin{align} \label{eq:state-sphere}
		\sqrt{ \frac{2 \pi^{d/2}}{\Gamma(d/2)} }\int_{\sS^{d-1}} | \vv \rangle {\rm d} \vv.
	\end{align}
\end{proof}

% \subsection{Proof of Proposition \ref{prop:Q-FO-estimator}}

\propqfoest*
To invoke Proposition \ref{prop:QVR}, we define the random variable $Z_x$ as
\begin{align*}
	Z_x = \nabla F (\vx+ \delta \vu; \xi), \quad {\rm where} \quad \vu \sim {\rm Unif}(\sB^d), ~~ \xi \sim \Xi.
\end{align*}
And in the proof of Proposition \ref{prop:FO-estimator}, we have shown that $ \E[Z_x] = \nabla f_{\delta}(\vx)$ and ${\rm Var}[Z_x] \le L^2$. 
Then by Proposition~\ref{prop:QVR} the \texttt{QuantumMeanEstimation+} can output a $(\delta,\sigma)$-SFO oracle with $\tilde \gO( \sqrt{d} L / \sigma)$ quantum sampling oracle of the random variable $Z_x$. Finally, we finish the proof by explicitly constructing the quantum sampling oracle, as we have done in the proof of Proposition \ref{prop:Q-ZO-estimator}. It is clear that it can be implemented by QSFO oracle (Definition \ref{dfn:QSFO}) given we have access to the quantum sampling oracle to $\vu$ and $\xi$. The quantum sampling oracle to $\xi \sim \Xi$ is just the same, and we only need to focus on generating the quantum sampling oracle to $\vu \sim {\rm Unif}(\sB^d)$. We achieve this by
applying the transform $| \vz \rangle \otimes | r \rangle \mapsto | r \vz \rangle \otimes | r \rangle $ to the uniform distribution over the unit sphere in (constructed as the same way we get Eq. \ref{eq:state-sphere}) and the uniform distribution over $[0,1]$ (constructed via Hadamard transform), which yields that
\begin{align*}
	\sqrt{ \frac{2 \pi^{d/2}}{\Gamma(d/2)} }\int_{\sS^{d-1}} | \vv \rangle {\rm d} \vv \otimes \int_{0}^1 | r \rangle {\rm d} r \mapsto \sqrt{\frac{\pi^{d/2}}{\Gamma(d/2+1)}} \int_{\sB^d} | \vu \rangle {\rm d} \vu \otimes \int_{0}^1 | r \rangle {\rm d} r.
\end{align*}

\section{Proof of the Complexity of Quantum Algorithms}

\thmQZO*
\begin{proof}
	Combining Theorem \ref{thm:o2ng} and Proposition \ref{prop:Q-ZO-estimator}, the total QSZO complexity to find a $(\delta,\epsilon)$-Goldstein stationary point is
	\begin{align*}
		\tilde \gO \left( \frac{\gamma + \delta L}{\delta}  \left( \frac{ d^{1/3} L^{2/3}}{\epsilon^{5/3}} + \frac{\sigma^2}{\epsilon^3} \right) \frac{dL}{\sigma}  \right).
	\end{align*}
	By selecting $\sigma= d^{1/6} L^{1/3} \epsilon^{2/3} $ to minimize the complexity above, we obtain the total QSZO complexity of $\tilde \gO(d^{7/6}(\gamma + \delta L) L^{4/3} \delta^{-1} \epsilon^{-7/3})$.
\end{proof}

\thmQFO*
\begin{proof}
	Combining Theorem \ref{thm:o2ng} and Proposition \ref{prop:Q-FO-estimator}, the total QSFO complexity to find a $(\delta,\epsilon)$-Goldstein stationary point is
	\begin{align*}
		\tilde \gO \left( \frac{\gamma + \delta L}{\delta}  \left( \frac{ d^{1/3} L^{2/3}}{\epsilon^{5/3}} + \frac{\sigma^2}{\epsilon^3} \right) \frac{\sqrt{d}L}{\sigma}  \right).
	\end{align*}
	By selecting $\sigma= d^{1/6} L^{1/3} \epsilon^{2/3} $ to minimize the complexity above, we obtain the total QSFO complexity of $\tilde \gO(d^{2/3}(\gamma + \delta L) L^{4/3} \delta^{-1} \epsilon^{-7/3})$.
\end{proof}

\section{Numerical Experiments}\label{app:numerical-experiments}

\subsection{First-order experiments}\label{app:fo-experiments}

\paragraph{Objective and distributed setting.}
We use the finite-sum synthetic objective
\begin{equation}\label{eq:fo-synthetic-objective}
	F(\vx;\xi)=\max_{1\leq r\leq R}
	\sin(\mathbf{a}_{\xi,r}^{\top}\vx+b_{\xi,r})
	+\lambda\lVert\vx\rVert_1,
\end{equation}
with $d=100$, $n=4096$, $R=1$, $\lambda=10^{-3}$, smoothing radius
$\delta=0.1$, feature scale one, zero phases, and a $0.25$ common offset in
the first feature coordinate. The sinusoid makes the objective nonconvex and
the $\ell_1$ term makes it nonsmooth. Because $R=1$, this experiment does not
exercise additional nonsmoothness from switching between multiple sinusoidal
branches. All methods start from the origin. The data are partitioned into
eight disjoint shards and optimized by eight CPU processes over a complete
mixing graph. We count one synchronized oracle layer as one unit of
\emph{depth}, and all training SFO calls across workers as \emph{work}; the
high-precision evaluation calls described below are excluded from training
work. The experiment consequently studies algorithmic communication/work
accounting, rather than multi-machine wall-clock speedup.

\paragraph{Pilot calibration and frozen configurations.}
Pilot seeds $100$--$104$ are used for parameter selection and formal seeds
$0$--$19$ are used only for reporting. The two sets are disjoint. Within the
same formal seed, both methods share the problem instance, data partition, and
evaluation bank, while retaining independent method-specific random streams.
Within the searched grids, the frozen NOG-FO configuration is $M=2$, $\eta=1$, one
smoothing sample, a global data batch of eight, and 960 main rounds. Its
perturbation radius is $\delta/M=0.05$; two explicitly counted initialization
queries give a maximum accounted depth of 962. The frozen ME-DOL-FO
configuration uses epoch length six, radius multiplier 100, one data sample per
worker per layer, and 3840 rounds. Its Euclidean domain radius and learning
rate are
\begin{equation}
	r_{\mathrm{ME}}=\frac{100\delta}{4\cdot6\sqrt{8}}\simeq0.1473,
	\qquad
	\eta_{\mathrm{ME}}=r_{\mathrm{ME}}/\sqrt{6}\simeq0.0601.
\end{equation}
The pilot objective minimizes the aggregate log-distance of the NOG/ME-DOL
work ratio from one, subject to all five pilot seeds hitting and to a
nondecreasing NOG batch schedule as $\epsilon$ decreases. These are
pilot-calibrated, locally good configurations within the searched grids, not
certified globally optimal parameters.

The complete v4 schedule changes the NOG batch from eight to 16 at
$\epsilon=0.01025$. Such a change affects depth and work simultaneously.
Accordingly, the homogeneous scaling comparison in the paper includes only
$\epsilon\geq0.0105$, for which the global batch remains eight; neither the
$0.01025$ nor the $0.010$ batch-16 result is used. A future extension below
this boundary should be rerun with the batch fixed at eight if it is to be
combined with the present curve.

\paragraph{Stationarity proxy and confirmed first hit.}
Computing the exact distance from the origin to the Goldstein subdifferential
is intractable for this stochastic test problem. At every checkpoint we
therefore estimate the norm of the gradient of the randomly smoothed objective
using 256 smoothing samples and 512 data samples per smoothing sample
($131{,}072$ evaluation samples). A method-independent evaluation bank is
fixed for each formal seed. NOG is evaluated every two depth layers and
ME-DOL at every six-layer epoch. To reduce false hits caused by Monte Carlo
noise, a threshold is declared reached only when two consecutive checkpoints
satisfy $\widehat{\operatorname{stat}}(\vx)\leq\epsilon$; the first
checkpoint of that pair defines the reported first-hit depth and work. This
quantity is a high-precision empirical proxy for Goldstein stationarity, not
the exact Goldstein subdifferential distance.

\paragraph{Statistical protocol.}
One trajectory per method and formal seed is reused to locate first hits for
all 25 thresholds, so thresholds within a seed are correlated. Absolute values
below are means $\pm$ sample standard deviations over 20 seeds. Ratios are
formed per seed and then averaged; they therefore need not equal ratios of the
two absolute means. Figure~\ref{fig:fo-v4-fixed-batch} uses 2000 paired
bootstrap resamples for the ratio confidence intervals. All 1000 reported
method--threshold--seed combinations reach the corresponding threshold within
their prescribed budgets, so no capped or censored observations enter the
table.

\begin{table}[t]
\centering
\caption{Complete FO results in the homogeneous batch-8 regime. Depth ratio
is ME-DOL/NOG and work ratio is NOG/ME-DOL; ratios are per-seed paired means.}
\label{tab:fo-v4-fixed-batch-full}
\scriptsize
\resizebox{\textwidth}{!}{%
\begin{tabular}{@{}rrrrrrr@{}}
\toprule
$\epsilon$ & NOG depth & ME-DOL depth & Depth ratio & NOG work & ME-DOL work & Work ratio \\
\midrule
0.20000 & 135.8 $\pm$ 1.7 & 66.3 $\pm$ 1.3 & 0.488 & 1086.4 $\pm$ 13.6 & 530.4 $\pm$ 10.7 & 2.049 \\
0.18000 & 162.6 $\pm$ 1.6 & 82.5 $\pm$ 2.7 & 0.507 & 1300.8 $\pm$ 12.8 & 660.0 $\pm$ 21.3 & 1.973 \\
0.16000 & 187.2 $\pm$ 1.8 & 96.0 $\pm$ 0.0 & 0.513 & 1497.6 $\pm$ 14.1 & 768.0 $\pm$ 0.0 & 1.950 \\
0.14000 & 211.2 $\pm$ 2.0 & 111.0 $\pm$ 3.1 & 0.526 & 1689.6 $\pm$ 15.9 & 888.0 $\pm$ 24.6 & 1.904 \\
0.12000 & 235.6 $\pm$ 2.3 & 127.5 $\pm$ 2.7 & 0.541 & 1884.8 $\pm$ 18.4 & 1020.0 $\pm$ 21.3 & 1.848 \\
0.10000 & 260.5 $\pm$ 3.0 & 146.1 $\pm$ 2.9 & 0.561 & 2084.0 $\pm$ 23.7 & 1168.8 $\pm$ 23.5 & 1.784 \\
0.09000 & 274.1 $\pm$ 3.5 & 156.9 $\pm$ 2.2 & 0.572 & 2192.8 $\pm$ 28.2 & 1255.2 $\pm$ 17.6 & 1.747 \\
0.08000 & 288.6 $\pm$ 4.2 & 168.6 $\pm$ 1.8 & 0.584 & 2308.8 $\pm$ 33.3 & 1348.8 $\pm$ 14.8 & 1.712 \\
0.07000 & 304.9 $\pm$ 3.9 & 182.4 $\pm$ 3.6 & 0.598 & 2439.2 $\pm$ 30.9 & 1459.2 $\pm$ 28.7 & 1.672 \\
0.06000 & 323.3 $\pm$ 4.4 & 198.0 $\pm$ 4.4 & 0.613 & 2586.4 $\pm$ 35.3 & 1584.0 $\pm$ 34.8 & 1.633 \\
0.05000 & 345.7 $\pm$ 5.8 & 217.2 $\pm$ 5.4 & 0.628 & 2765.6 $\pm$ 46.2 & 1737.6 $\pm$ 42.9 & 1.593 \\
0.04000 & 372.4 $\pm$ 6.4 & 241.5 $\pm$ 6.7 & 0.649 & 2979.2 $\pm$ 51.5 & 1932.0 $\pm$ 53.7 & 1.543 \\
0.03000 & 408.4 $\pm$ 10.0 & 275.1 $\pm$ 13.8 & 0.674 & 3267.2 $\pm$ 80.2 & 2200.8 $\pm$ 110.4 & 1.489 \\
0.02500 & 429.8 $\pm$ 12.4 & 297.6 $\pm$ 15.7 & 0.693 & 3438.4 $\pm$ 99.4 & 2380.8 $\pm$ 125.9 & 1.448 \\
0.02000 & 459.9 $\pm$ 14.4 & 330.9 $\pm$ 20.6 & 0.720 & 3679.2 $\pm$ 115.4 & 2647.2 $\pm$ 165.0 & 1.395 \\
0.01800 & 476.6 $\pm$ 20.8 & 347.7 $\pm$ 22.9 & 0.731 & 3812.8 $\pm$ 166.2 & 2781.6 $\pm$ 183.3 & 1.376 \\
0.01600 & 492.4 $\pm$ 20.9 & 369.0 $\pm$ 25.3 & 0.751 & 3939.2 $\pm$ 167.5 & 2952.0 $\pm$ 202.2 & 1.340 \\
0.01500 & 499.2 $\pm$ 19.7 & 381.3 $\pm$ 25.1 & 0.765 & 3993.6 $\pm$ 157.4 & 3050.4 $\pm$ 200.9 & 1.315 \\
0.01400 & 508.0 $\pm$ 19.9 & 394.8 $\pm$ 28.1 & 0.778 & 4064.0 $\pm$ 159.2 & 3158.4 $\pm$ 224.9 & 1.293 \\
0.01300 & 518.9 $\pm$ 23.2 & 418.8 $\pm$ 35.6 & 0.809 & 4151.2 $\pm$ 185.6 & 3350.4 $\pm$ 284.9 & 1.248 \\
0.01200 & 538.9 $\pm$ 28.5 & 446.1 $\pm$ 51.0 & 0.831 & 4311.2 $\pm$ 227.7 & 3568.8 $\pm$ 408.0 & 1.222 \\
0.01150 & 546.0 $\pm$ 28.2 & 468.3 $\pm$ 61.5 & 0.861 & 4368.0 $\pm$ 225.9 & 3746.4 $\pm$ 491.6 & 1.183 \\
0.01100 & 566.2 $\pm$ 29.3 & 535.8 $\pm$ 107.7 & 0.948 & 4529.6 $\pm$ 234.7 & 4286.4 $\pm$ 861.8 & 1.091 \\
0.01075 & 574.1 $\pm$ 30.3 & 590.4 $\pm$ 133.5 & 1.036 & 4592.8 $\pm$ 242.7 & 4723.2 $\pm$ 1067.8 & 1.019 \\
0.01050 & 587.2 $\pm$ 37.7 & 642.9 $\pm$ 197.4 & 1.102 & 4697.6 $\pm$ 301.2 & 5143.2 $\pm$ 1579.4 & 0.980 \\
\bottomrule
\end{tabular}}
\end{table}

\paragraph{Interpretation relative to Theorem~\ref{thm:parallel-FO}.}
For first-order oracles, the theorem gives NOG depth
$\widetilde O(d^{1/3}\delta^{-1}\epsilon^{-5/3})$ and work
$\widetilde O(\delta^{-1}\epsilon^{-3})$, whereas ME-DOL has the reference
depth and work dependence $\widetilde O(\delta^{-1}\epsilon^{-3})$. Thus the
worst-case bounds suggest a relative depth factor proportional to
$\epsilon^{-4/3}/d^{1/3}$ while the two work complexities remain of the same
order. On the fixed-batch grid, the observed paired depth ratio rises
monotonically from 0.488 to 1.102, while the paired work ratio stays between
0.980 and 2.049. The latter maximum slightly exceeds the preregistered
matched-work window $[0.5,2]$ at $\epsilon=0.2$, so we describe the work as
order-one rather than claiming that every preregistered criterion passes.

The depth and work ratios in this particular fixed-batch design are also not
independent measurements: both methods spend eight SFO calls per depth layer,
so, up to NOG's explicitly counted initialization, their per-seed ratios are
reciprocally coupled. Moreover, worst-case upper bounds need not be tight on
one finite synthetic instance, and the present $\epsilon$ range need not be
asymptotic. We therefore use the experiment only as qualitative evidence for
the predicted improvement in relative depth without an order-of-magnitude work
gap; it neither proves the theory nor identifies the exact $5/3$, $3$, or
$4/3$ exponents.

\paragraph{Reproducibility and audit.}
The supplementary artifact stores the full configuration, frozen schedule,
per-seed summaries, paired ratios, and SHA256 audit records under
\path{results/theory_validation_v4}. The formal audit checks all 60 underlying
trajectories for seed/configuration identity, strictly increasing depth and
work, exact total and per-worker SFO accounting, and finite nonnegative
stationarity proxies; all 60 pass. From the repository root, the complete
pipeline is reproduced by successively running the modules
\path{src.distributed.theory_validation_runner} (with
\path{pilot-batch-grid}, then \path{formal}),
\path{src.distributed.theory_validation_freeze},
\path{src.distributed.theory_validation_audit},
\path{src.distributed.theory_validation_analysis},
\path{src.distributed.theory_validation_report}, and
\path{src.distributed.theory_validation_package} in the \path{NOG} Conda
environment.


% \section{Proof of the Lower Bound}

% We recall the (unscaled) hard instance for smooth nonconvex optimization, introduced by \citet{carmon2020lower}:
% \begin{align} \label{eq:hard-nc}
	%    \bar f_T(\vx) = - \Psi(1) \Phi(x_1) + \sum_{i=2}^T [ \Psi(-x_{i-1} \Phi(-x_i) - \Psi(x_{i-1}) \Phi(x_i)],
	% \end{align}
% where
% \begin{align*}
	%     \Psi(x) &= \begin{cases}
		%         0, & x \le 1/2, \\
		%         \exp(1 - 1/(2x-1)^2), & x > 1/2,
		%     \end{cases} \quad {\rm and} \quad
	%     \Phi(x) = \sqrt{\rm e} \int_{-\infty}^x \exp (-t^2/2) {\rm d} t.
	% \end{align*}
% We list the properties of this function as follows. They can all be found in different parts of \citep{carmon2020lower}, and they are summarized in \citep[Lemma 2]{arjevani2023lower} with explicit numerical constants. We restate them as follows. 
% \begin{lem}[{\citep[Lemma 2]{arjevani2023lower}}] \label{lem:Carmon}
	% Let ${\rm prog}_c (\vx) = \max\{i : \vert x_i \vert \ge c  \}$. The hard instance defined in (\ref{eq:hard-nc}) satisfies 
	% \begin{enumerate}
		%     \item[(a)] $\bar f_T(\vzero) = 0 $ and $\inf_{\vx \in \sR^d} \bar f_T(\vx) \ge - \gamma_0 T$ with $\gamma_0 = 12$.
		%     \item[(b)] $\bar f_T(\,\cdot\,)$ is $H_0$-smooth with $H_0 = 152$. 
		%     \item[(c)] $\bar f_T(\,\cdot\,)$ is $L_0 \sqrt{T}$-Lipschitz with $L_0 = 23$.
		%     \item[(d)] ${\rm prog}_0  (\nabla \bar f_T(\vx)) \leq {\rm prog}_{1/2} (\vx) + 1$ 
		%     \item[(e)] ${\rm prog}_1(\vx) < T \Rightarrow \Vert \nabla \bar f_T(\vx) \Vert \ge \vert \nabla \bar f_T(\vx)_{{\rm prog}_1(\vx)+1} \vert \ge 1$.
		% \end{enumerate}
	% \end{lem}

% Lemma \ref{lem:Carmon}(d) means that any \textit{zero-respecting} algorithms defined in \citep{carmon2020lower}, which are the algorithms that never explore directions with zero derivatives, can only make at least one more zero coordinate become non-zero with one first-order oracle call.
% In this way, the progress of any zero-respecting algorithm is restricted to a chain-like structure, and Lemma \ref{lem:Carmon}(e) indicates that the algorithm needs to make $\Omega(T)$ first-order oracle calls to reach the end of the chain in order to find a point with small gradient norm. 
% The above properties lead to an $\Omega(\epsilon^{-2})$ lower bound in the classical setting by appropriately scaling the function: simple calculus shows that the largest possible $T$ to meet the regular conditions for the objective and ensure a lower bound for finding $\epsilon$-stationary points.

% As shown in \citep{carmon2020lower},
% the lower bound for zero-respecting algorithms can further be extended to any deterministic algorithms using reduction via \textit{resisting oracles} \citep{nemirovskij1983problem}, and to any randomized algorithms by defining a distribution over the rotated function % Combined with Lemma \ref{lem:Carmon}(e), any zero-respecting algorithm beginning at $\vzero$ requires $\Omega(T)$ first-order oracle calls to reach a $1$-stationary point. Then, by appropriately scaling the function, we can conclude an $\Omega(\epsilon^{-2})$  
% % require at least $\Omega(\epsilon^{-2})$ first-order oracle complexity to find an $\epsilon$-stationary point of $\bar f_T(\vx)$ \citep[Theorem 1]{carmon2020lower}.
% % \citet{carmon2020lower} proved that the function defined in (\ref{eq:hard-nc}) is hard for any \textit{zero-respecting} algorithms such that they require at least $\Omega(\epsilon^{-2})$ first-order queries to find an $\epsilon$-stationary point of $\nabla \bar f_T(\,\cdot\,)$. Furthermore, for any classical randomized algorithm \texttt{A} which produces iterates of the form
% % \begin{align*}
	% %     \vx_t = \texttt{A}(r, \nabla f(\vx_1), \cdots, \nabla f(\vx_{t-1})),
	% % \end{align*}
% % where $r $ is a random uniform variable on $[0,1]$, one can also prove an $\Omega(\epsilon^{-2})$ lower bound by define the rotated function
% \begin{align} \label{eq:rotated-hatd-instance}
	%     \bar f_{T;\mU}(\vx): = \bar f_{T}(\mU^T \vx),
	% \end{align}
% where $\mU \in \sR^{d \times T}$ (with columns $ \vu_1,\cdots,\vu_T $) is an orthogonal matrix chosen at random from the space of orthogonal matrices ${\rm Orth}(d,T) = \{ \mV \in \sR^{d \times T} \mid \mV^\top \mV = \mI_T \}$. 

% The lower bound under the stochastic setting also relies on the above hard instance under the deterministic setting. Recently, \citet{zhang2023quantum} embedded the lower bound for quantum mean estimation under the stochastic setting to the probability of each coordinate being activated (changing from non-zero to zero) \textcolor{red}{is $\gO(\epsilon^2)$}. This demonstrates an $\Omega(\epsilon^{-4})$ lower bound for quantum stochastic-first-order algorithm. 

% % $\mU \in \sR^{d \times T}$ is a random orthogonal matrix that is chosen uniformly from the space of orthogonal matrices $ \{ \mU \in \sR^{d \times T}: \mU^\top \mU = \mI_{T} \}$. In the following, we also let $\vu_1,\cdots,\vu_T$ to be the columns of $\mU$ and define $\mU_t$ be the sub-matrix generated by the first $t$ columns of $\mU$.

% Given the above facts, \citet{zhang2023quantum} explicitly construct a worst-case stochastic gradient estimator $\bar \vg_{T; \mU}(\vx;\xi)$ for any $\bar f_{T; \mU}$ that satisfies Assumption \ref{asm:FO}(b), showing that an $\Omega(\epsilon^{-4})$ lower bound for any quantum algorithms in cases $d = \Omega(\epsilon^{-4})$, which generalizes the lower bound for stochastic smooth nonconvex optimization in the classical setting \citep{arjevani2023lower}. Specifically, \citet{zhang2023quantum} constructed the following QS





% \begin{algorithm}[t] \label{alg:hard-QSFO}
	%     \caption{The Quantum Stochastic First-Order Oracle}
	%     \begin{algorithmic}[1]
		%     {\Procedure{\texttt{Initialize}}{$ $ }
			%     \State Sample $\mU \in \sR^{d \times T}$ uniformly from ${\rm Orth}(d,T)$.
			%     \State Create $T$ disjoint orthogonal subspaces $ \fV_1,\cdots,  \fV_T$ from ${\rm Orth}(d,T)$.  
			%     % \While{$x>-5$}
			%     % \State $x\gets x-1$ 
			%     % \EndWhile 
			%     \EndProcedure   }
		%     {\Procedure{\texttt{Query}}{$\vx$}
			%     \State 
			%     \EndProcedure
			%     }
		%     \end{algorithmic}
	%     \end{algorithm} 

% \citet{zhang2023quantum} uses the hybrid argument (which is developed by \citet{garg2020no}) to prove their quantum lower bound. Specifically, for any quantum algorithm $\bar \mA_{\rm quan}$ that queries no more than $KS/2$ calls of oracles 
% \begin{align*}
	%     \bar \mO_{T;\mU}: | \vx \rangle \otimes | \xi \rangle \otimes | \vy \rangle \mapsto | \vx \rangle \otimes | \xi \rangle \otimes | \bar \vg_{T; \mU}(\vx, \xi) + \vy \rangle,
	% \end{align*}
% we can define a sequence of unitary transforms as follows (where $\mA_0 = \bar \mA_{\rm quan})$:
% \begin{align} \label{eq:hybrid}
	% \begin{split}
		%         \mA_0 &:= \mB_{K+1} \mB_{K;\mU}  \cdots \mB_{2;\mU} \mB_{1;\mU} \\
		%     \mA_1 &:= \mB_{K+1} \mB_{K; \mU}  \cdots \mB_{2; \mU}  \mB_{1;\mU_1} \\
		%     \mA_2 &:= \mB_{K+1} \mB_{K; \mU} \cdots \mB_{2; \mU_2}  \mB_{1;\mU_1} \\
		%     \vdots & \\
		%     \mA_K &:= \mB_{K+1} \mB_{K; \mU_K} \cdots \mB_{2; \mU_2}  \mB_{1;\mU_1},
		% \end{split}
	% \end{align}
% where $\mB_{k;\mU_k} $ takes the form of
% \begin{align*}
	%     \mB_{k; \mU_k} := \bar \mO_{k:\mU_k} \mV_{k, S/2} \cdots \bar \mO_{k; \mU_k} \mV_{k,2} \bar \mO_{k, \mU_k} \mV_{k,1}.
	% \end{align*}
% In the above $\mV_{k,1},\cdots ,\mV_{k,S/2}$ are unitary transforms by the basic postulate in quantum mechanisms \citep[Postulate 2]{nielsen2010quantum}, which comes from the celebrated Schrodinger equation. For the above hybrids, \citep[Lemma D.3]{zhang2023quantum} claims the following indistinguishable result if all the iterates lie in a bounded set.

% % \begin{asm} \label{asm:bounded}
	% %  Assume the quantum algorithm $\bar \mA_{\rm quan}$ only queries stochastic gradients of the points $\vx$ that lie in the bounded set $\sB(\vzero, 2 \sqrt{T} )$.    
	% % \end{asm}

% \begin{lem}[{\citep[Lemma D.3 and D.10]{zhang2023quantum}}] \label{lem:chenyi}
	% Recall the numerical constant defined in Lemma \ref{lem:Carmon}(c) that $L_0= 23$.
	% Let $d \ge 2 S T \log S$, then there exists a distribution of functions $f_{T;\mU}$ defined in (\ref{eq:rotated-hatd-instance}) and a distribution over stochastic gradients $\bar \vg_{T; \mU}$ that satisfies $\E_{\xi}[\bar \vg_{T;\mU}(\vx; \xi)] = \nabla f_{T;\mU}(\vx)$ and ${\rm Var}_{\xi}[\bar \vg_{T;\mU}] \le 4 L_0^2 S$, such that the hybrids defined in (\ref{eq:hybrid}) with $ \mA_0 = \bar \mA_{\rm quan}$ satisfies
	% \begin{align*}
		%     \E_{\vu_t,\cdots, \vu_T, \zeta} \Vert \mA_t | \vzero \rangle - \mA_{t-1} | \vzero \rangle \Vert^2 \le {1}/{(72 T^4)},
		% \end{align*}
	% where $\zeta$ represents the randomness in generating $\bar \vg_{T;\mU}$, and $\mU$ does not depend on $\zeta$.
	% \end{lem}

% By apply Lemma \ref{lem:chenyi} recursively on the hybrids defined in (\ref{eq:hybrid}), we can conclude that $\mA_{\rm quan}$ will have outputs similar to $\mA_K$. \citet{zhang2023quantum} used the above lemma to prove their lower bound, by arguing that $\mA_K$ can not find an $\epsilon$-stationary point because conditional  $\mU_K$ the vector $\vu_T$ still looks random for $\mA_K$ when $K < T$.
% However, the proof of \citet{zhang2023quantum} is specified for smooth objectives, and in the following we adapt their analysis to our setting by using a similar technique in \citep[Section 7]{arjevani2023lower}.
% Inspired by \citet{cutkosky2023optimal},  we define the following auxiliary functions
% \begin{align} \label{eq:func-aux}
	%     \rho_R(\vx) = \frac{\vx}{\sqrt{1 + \Vert \vx \Vert^2 / R^2}}, \quad q_R(\vx) =  \vx^\top \rho_R(\vx).
	% \end{align}
% And then define the regularized function as
% \begin{align} \label{eq:reg-hard-ins}
	%     \hat f_{T;\mU}(\vx) = \bar f_{T; \mU}(\rho_R(\vx)) + \lambda q_R(\vx).
	% \end{align}
% In the above, the bijective function $\rho_R(\vx)$ projects $\vx$ onto the bounded set, and the Lipschitz-type regularizer $q_R(\vx)$ penalizes the points with large norm such that they would have large gradients. 
% The regularized objective in (\ref{eq:reg-hard-ins}) has the following properties.
% \begin{lem}[{\citep{cutkosky2023optimal}}] \label{lem:ashok-25}
	% Let $\lambda = 1/10$ and $R = 10 L_0 \sqrt{T}$ in (\ref{eq:reg-hard-ins}), where $L_0=23$ is the numerical constant defined in Lemma \ref{lem:Carmon}(c). Then the function $\hat f_{T; \mU}(\vx)$ satisfies
	% \begin{enumerate}
		%  \item[(a)] $\hat f_{T;\mU}(\vzero) = 0 $ and $\inf_{\vx \in \sR^d} \bar f_{T;\mU}(\vx) \ge - \gamma_0 T$ with $\gamma_0 = 12$.
		%     \item[(b)] $\hat f_{T;\mU}(\,\cdot\,)$ is $\hat H_0$-smooth with $\hat H_0 = 156$. 
		%     \item[(c)] $\hat f_{T;\mU}(\,\cdot\,)$ is $\hat L_0 \sqrt{T}$-Lipschitz with $\hat L_0 = 92$.
		%     \item[(d)] Let $\vy = \rho_R(\vx)$. If ${\rm prog}_{1/4} (\mU^\top \vy) <T$, then $\Vert \nabla \hat f_{T; \mU}(\vx) \Vert \ge 1/2$.   
		% \end{enumerate} % \begin{align*}
		% % \end{align*}
	% \end{lem}

% We remark that the above (a-c) can be found in {\citep[Lemma 26]{cutkosky2023optimal}}, while (d) appears in {\citep[Lemma 25]{cutkosky2023optimal}}.
% For this hard instance $\hat f_{T; \mU}$,
% We further define the stochastic gradient associated as
% \begin{align} \label{eq:stoc-grad-reg}
	%     \hat \vg_{T; \mU}(\vx;\xi) = \mJ_{\rho_R}(\vx)^\top \bar \vg_{T; \mU}(\rho_R(\vx);\xi) + \lambda \nabla q_R(\vx),
	% \end{align}
% where $\mJ_{\rho_R}(\vx)$ denotes the Jacobian of the function $\rho_R$ evaluated at the point $\vx$, which by calculation is
% \begin{align} \label{eq:Jacobian}
	%     \mJ_{\rho_R}(\vx) = \frac{\mI_d - \rho_R(\vx) \rho_R(\vx)^\top / R^2}{\sqrt{1 + \Vert \vx \Vert^2 / R^2}}.
	% \end{align}
% The above stochastic gradient estimator (\ref{eq:reg-hard-ins}) induces
% the following QSFO oracle 
% \begin{align} \label{eq:QSFO-reg}
	%     \hat \mO_{T;\mU}: | \vx \rangle \otimes | \xi \rangle \otimes | \vy \rangle \mapsto | \vx \rangle \otimes | \xi \rangle \otimes | \hat \vg_{T; \mU}(\vx, \xi) + \vy \rangle.
	% \end{align}
% % Now we define $\vz = \rho_R(\vx)$ and recall \citep[Lemma 25]{cutkosky2023optimal}.

% % \begin{lem}[{\citep[Lemma 25]{cutkosky2023optimal}}] \label{lem:ashok-25}
	% % Recall the numerical constant defined in Lemma \ref{lem:Carmon}(c) that $L_0= 23$. If we define $R = 10 L_0 \sqrt{T}$ and $\lambda = 1/10$ in (\ref{eq:reg-hard-ins}), then
	% % \begin{align*}
		% %     {\rm prog}_{1/4}(\mU^\top \vz) <T \quad \Rightarrow \quad \Vert \nabla \hat f_{T; \mU}(\vx) \Vert \ge 1/2.
		% % \end{align*}
	% % \end{lem}

% Below, we formally prove the following key result.

% \begin{thm} \label{thm:hard-T2}
	% Under the same parameter setting as Lemma \ref{lem:chenyi} and Lemma \ref{lem:ashok-25}
	% there exist a distribution of function $\hat f_{T;\mU} $ defined as (\ref{eq:reg-hard-ins}) and a distribution of QSFO defined as (\ref{eq:QSFO-reg}), such that 
	% any quantum algorithm $\hat \mA_{\rm quan}$ that uses less than $(T S/ 2)$ randomly selected QSFO calls cannot output an $(1/2)$-stationary point of $ \hat f_{T; \mU}$ with probability larger than $1/2$. Specifically, let $ p_{\rm out} $ be the probability distribution over $\vx $ obtained by measuring the state $ \hat \mA_{\rm quan} | \vzero \rangle $, we have
	% \begin{align*}
		%     \BP_{\vx \sim p_{\rm out}}( \Vert \nabla \hat f_{T; \mU}(\vx) \Vert \le 1/2  ) < 1/2.
		% \end{align*}
	% % then there exists a randomized wors-case stochastic gradient $\bar \vg_{T; \mU}$ of the rotated function $\bar f_{T;\mU}$ defined in (\ref{eq:rotated-hatd-instance}) that satisfies Assumption \ref{asm:FO}(b), such that
	% %     Let xxx. Let xxx. There exists a  such that  $\gO(\delta^{-1} \epsilon^{-3})$ QSZO calls can not xxx.
	% \end{thm}

% \begin{proof}
	% Let $\vz = \rho_R(\vx)$. It is clear that there is a bijection between the variables $\vx$ and $\vz$. Moreover, we know there also exist a bijection between the stochastic gradient estimators $  \hat \vg_{T; \mU}(\vx; \xi) $ and $\bar \vg_{T;\mU}(\vz; \xi)$: Given $\vz$ and $\bar \vg_{T;\mU}(\vz,\xi)$ we can calculate $\vx = \rho_R^{-1}(\vz)$ and then use the formula in (\ref{eq:QSFO-reg}) to calculate $\hat \vg_{T;\mU}(\vx;\xi)$; Conversely, given $\vx$ and $\hat \vg_{T;\mU}(\vx;\xi)$, we can also obtain $\bar \vg_{T;\mU}(\vz,\xi)$ from $\bar \vg_{T;\mU}(\vz,\xi) = [\mJ_{\rho_R}]^{-1} (\hat \vg_{T;\mU}(\vx;\xi) - \lambda \nabla q_R(\vx))$. Therefore, for the quantum algorithm $\hat \mA_{\rm quan}$ running on $\vx$ which uses $N$ oracle  calls of $\hat \mO_{T;\mU}$, there exists an one-to-one mapping of another quantum algorithm $\bar \mA_{\rm quan}$ running on $\vz = \rho_R(\vx)$ which uses $N$ oracle calls of $\bar \mO(T;\mU)$. 
	% Next, we invoke Lemma \ref{lem:chenyi} recursively to get 
	% \begin{align*}
		%     \E_{\mU, \zeta} \Vert \mA_K | \vzero \rangle - \mA_0 | \vzero \rangle \Vert^2 \le K \sum_{t=1}^{K-1} \E_{\mU, \zeta} \Vert \mA_{t+1} | \vzero \rangle - \mA_t | \vzero \rangle \Vert^2 \le \frac{1}{72T^2}.
		% \end{align*}
	% By Markov's inequality, with probability at least $5/6$ it holds that
	% % there exists an event $F$ with $\BP(F) \ge 2/3$ such that
	% \begin{align} \label{eq:norm-Ak-A0}
		%     \Vert \mA_K | \vzero \rangle - \mA_0 | \vzero \rangle \Vert^2 \le \frac{1}{60 T^2}.
		% \end{align}
	% We let $F$ be the event that the above happens, so $\BP(F) \ge 5/6$.
	
	% Recall the definition of $\mA_K$ in (\ref{eq:hybrid}), we know that the output of $\mA_K$ is fully determined by $\{ \vu_1,\cdots,\vu_K\}$ and the random variable $\zeta$ which represents the randomness in the QSFO calls. Note that 
	% in the construction of Lemma \ref{lem:chenyi} $\mU$ does not depend on $\zeta$. Therefore, conditioned on $\{ \vu_1,\cdots,\vu_K\}$ and $\zeta$, $\vu_T$ is a random unit vector in the orthogonal complement of the subspace spanned by $\{ \vu_1,\cdots,\vu_K\}$. In other words, $\vu_T$ follows a uniform distribution over the unit sphere in a $(d-K)$ space. By a standard concentration of measure
	% bound on the sphere (see, \textit{e.g.} \citep[Lemma 2.2]{ball1997elementary}), for any vector $\vy \in \sR^d$ which is independent of $ \vu_T$ it holds 
	% \begin{align*}
		%     \BP( \vert \langle \vu_T, \vy \rangle  > c \Vert \vy \Vert ) \le 2 \exp( - (d-K) c^2 /2).
		% \end{align*}
	% Plugging into $c = 1/ (4R)$ in the above and recall in Lemma \ref{lem:ashok-25} that $R = 10 L_0 \sqrt{T}$, we get 
	% \begin{align*}
		%     \BP( \vert \langle \vu_T, \vy \rangle \vert > \Vert \vy \Vert/ (4R) ) \le 2 \exp( -(d-K) / (3600T) ) .
		% \end{align*}
	% Taking $d = \Omega(T)$, specifically, $d \ge (3600 \ln 12 +1)T$, then with probability at least $5/6$, it holds that 
	% \begin{align} \label{eq:bound-uT}
		%     \vert \langle \vu_T, \vy \rangle \vert \le 1/4, \quad \forall \vy \in \sR(\vzero, R) \text{ that is independent of } \vu_T.
		% \end{align} % Now we let $E$ be the event that the above happens, so $\BP(E) \ge 5/6$. Note that $G$ is taken to be independent of $\mA_K$. 
	% % As we have discussed, any output of the quantum algorithm $\mA_K$ are independent of $\vu_T$. 
	% Recall the event $F$ defined as (\ref{eq:norm-Ak-A0}) happens and $\BP(F) \ge 5/6$. We know that 
	% \begin{align*}
		%     \BP(E \cap F) = 1 - \BP(E^c \cup F^c) \ge 1 - \BP(E^c) - \BP(F^c) \ge 2/3.  
		% \end{align*}
	% We let $G = E \cap F$.
	% Now, define the sets of $\vz = \rho_R(\vx)$ such that $\vx$ is a $(1/2)$-stationary point: $S =  \{ \vz \in \sR^d : \Vert \nabla \hat f_{T;\mU}(\rho_R^{-1}(\vz)) \Vert \le 1/2  \}$. 
	% We let $\psi_{K \mid G}(\vz)$ and $\psi_{0 \mid G}(\vz)$ to represent the wave functions associated with the states $\mA_K | \vzero \rangle$ and $\mA_0 | \vzero \rangle$ conditioned on the event $G$. 
	% Then 
	% $ p_{K \mid G} = \vert \psi_{K \mid G} \vert^2$ and $ p_{0 \mid G} = \vert \psi_{0 \mid G} \vert^2 $ represent the 
	% probability distribution obtained by measuring the states $\mA_K | \vzero \rangle$ and $\mA_0 | \vzero \rangle$ conditioned on the event $G$. 
	% The definition of total variation distance indicates
	% \begin{align} \label{eq:P-TV}
		%     \BP( \vz \in S \mid \vz \sim p_{0 \mid G}, G ) \le  \underbrace{\BP( \vz \in S \mid \vz \sim p_{K \mid G}, G  )}_{\rm (I)} + \underbrace{d_{\rm TV}(p_{0 \mid G}, p_{K \mid G})}_{\rm (II)}.
		% \end{align}
	% Note that by Lemma \ref{lem:ashok-25}(a) and (\ref{eq:bound-uT}), (I) above is always zero. Therefore we only need to upper bound (II). By the equivalent definition of total variation distance \citep[Proposition 4.2]{levin2017markov}, we have
	% \begin{align*}
		% d_{\rm TV} (p_{0 \mid G}, p_{K \mid G})  
		%      &= \frac{1}{2}\int \left\vert p_{0 \mid G}(\vz) - p_{K \mid G}(\vz) \right \vert {\rm d} \vz \\
		%     &= \frac{1}{2}\int \left(\vert \psi_{0 \mid G}(\vz) \vert + \vert \psi_{K \mid G}(\vz) \vert \right) \cdot \left \vert \psi_{0 \mid G}(\vz) - \psi_{K \mid G}(\vz) \right \vert {\rm d} \vz \\
		%     &\le \frac{1}{2} \sqrt{ \int  \left(\vert \psi_{0 \mid G}(\vz) \vert + \vert \psi_{K \mid G}(\vz) \vert \right)^2 {\rm d} \vz} \sqrt{ \int  \left \vert \psi_{0 \mid G}(\vz) - \psi_{K \mid G}(\vz) \right \vert^2 {\rm d} \vz} \\
		%     &\le \sqrt{ \int  \left \vert \psi_{0 \mid G}(\vz) - \psi_{ K \mid G}(\vz) \right \vert^2 {\rm d} \vz}  \le \frac{1}{2 \sqrt{15} T}.
		% \end{align*}
	% Above, the second last is from the normalized condition $\int \vert \psi_{0 \mid G}(\vz) \vert^2 {\rm d} \vz = 1$ and $\int \vert \psi_{K \mid G}(\vz) \vert^2 {\rm d} \vz = 1$, and the last inequality is due to (\ref{eq:norm-Ak-A0}). Because $T \ge 1$, (\ref{eq:P-TV}) becomes
	% \begin{align*}
		%     \BP( \vz \in S \mid \vz \sim p_{0 \mid G}, G ) \le  \frac{1}{2 \sqrt{15}}.
		% \end{align*}
	% Then we can conclude
	% \begin{align*}
		%     \BP( \vz \in S \mid \vz \sim p_{0}) \le \BP(\vz \in S \mid \vz \sim p_{0 \mid G}, G) + \BP(G^c) < \frac{1}{2}.
		% \end{align*}
	
	% \end{proof}
% Finally, we scale the hard instance above to prove the lower bound for quantum algorithms in finding $(\delta,\epsilon)$-stationary points.
% \thmlb* 

% \begin{proof}
	%    We scale the hard instance in Theorem \ref{thm:hard-T2} by $f_{T; \mU}(\vx) = \alpha \hat f_{T,\mU}(\vx / \beta)$ to simultaneously fulfill all the assumptions on the function class, as detailed below.
	%    \paragraph{I. Can not Find Stationary Points} Note that $\nabla f_{T; \mU} (\vx) = (\alpha / \beta) \nabla \hat f_{T; \mU}(\vx/ \beta)$, Theorem \ref{thm:hard-T2} that any quantum algorithm can not find an $(1/2)$-stationary point of $\hat f_{T; \mU}(\vx)$ implies that any quantum algorithm can not find an $(\alpha/(2 \beta))$-stationary point of $f_{T; \mU}(\vx)$. We let $\alpha = 6 \epsilon \beta $ to give a lower bound for finding $3\epsilon$-stationary points. 
	%    \paragraph{II. Can not Find Goldstein Stationary Points} Our next goal is to make the lower bound for finding $\epsilon$-stationary points to the lower bound for finding $(\delta,\epsilon)$-stationary points. This is where Proposition \ref{prop:zhang-rela} comes into play. From Proposition \ref{prop:zhang-rela}  we know that for an $H$-smooth objective, the lower bound for finding $3 \epsilon$-stationary point also implies the lower bound for finding $(\delta,\epsilon)$-stationary points with $\delta \le \epsilon / H$. And the smooth constant can be calculated as follows: By Lemma \ref{lem:ashok-25}(b) the unscaled function $\hat f_{T;\mU}$ is $\hat H_0$-smooth with $\hat H_0=156$, so the scaled function is $(\alpha \hat H_0/ \beta^2) $-smooth. Therefore, we can set $\alpha \hat H_0  / \beta^2 = \epsilon / \delta$.
	%    \paragraph{III. Bounded Initial Suboptimality Gap} 
	%    By Lemma \ref{lem:ashok-25}(a), we have $f_{T,\mU}(\vzero) = 0$ and $\inf_{\vx \in \sR^d} f(\vx) \ge - \gamma_0 \alpha T$, where $\gamma_0  =12$. Therefore, we can set $T = \gamma / (\gamma_0 \alpha)$ to fulfill the condition $f_{T; \mU}(\vzero) - \inf_{\vx \in \sR^d} f(\vx) \le \gamma$. 
	%     \paragraph{IV. Bounded Variance}  From $\nabla f_{T; \mU} (\vx) = (\alpha / \beta) \nabla \hat f_{T; \mU}(\vx/ \beta)$, the stochastic gradient estimator is also scaled as $ \vg_{T; \mU}(\vx; \xi) = (\alpha/ \beta) \hat \vg_{T; \mU}(\vx / \beta; \xi) $. Recall the definition of $\hat \vg_{T;\mU}$ in (\ref{eq:stoc-grad-reg}), then we have
	%    \begin{align*}
		%         \vg_{T; \mU}(\vx; \xi)  = (\alpha/ \beta) \left( \mJ_{\rho_R}(\vx / \beta)^\top \bar \vg_{T; \mU}(\rho_R(\vx / \beta);\xi) + \lambda \nabla q_R(\vx/ \beta) \right). 
		%    \end{align*}
	%    From (\ref{eq:Jacobian}) we know that $\Vert \mJ_{\rho_R}(\,\cdot\,) \Vert \le 1$, and by Lemma \ref{lem:chenyi} we know ${\rm Var}(\bar \vg_{T; \mU}) \le 2 L_0 \sqrt{S}$, where $L_0 = 23$, so we can conclude that $ {\rm Var}(\vg_{T;\mU}) \le (2 \alpha L_0 \sqrt{S} / \beta)^2$. Therefore, we can let $\beta \ge 2 \alpha L_0 \sqrt{2S} / L$ to fulfill the condition that ${\rm Var}(\vg_{T;\mU}) \le L^2 /2$. 
	%    \paragraph{V. Bounded Gradient} Note that $\nabla f_{T; \mU} (\vx) = (\alpha / \beta) \nabla \hat f_{T; \mU}(\vx/ \beta)$. Recall Lemma \ref{lem:ashok-25}(c) that $\hat f_{T; \mU}(\vx)$ is $\hat L_0 \sqrt{T}$-Lipschitz, we have that the scaled function $f_{T;\mU}(\vx)$ is $(\alpha \hat L_0 \sqrt{T}/ \beta)$-Lipschitz, where $\hat L_0 = 92$. Therefore, we can let $\beta \ge \alpha \hat L_0  \sqrt{2T}/ L$ to fulfill the condition that $f_{T; \mU}(\vx)$ is $(L/\sqrt{2})$-Lipschitz. 
	%    \paragraph{VI. Final Lower Bound} 
	%    We specify the parameter setting from the above analysis. From \textbf{I} we have $\alpha = 6 \epsilon \beta$. Substitute it into \textbf{II} yields that $\beta = 6 \hat H_0 \delta$. Hence we obtain $\alpha = 36 \hat H_0 \epsilon \delta$. Then from \textbf{III} we have 
	%    \begin{align*}
		%        T = \frac{\gamma}{\gamma_0 \alpha} = \frac{\gamma}{36 \gamma_0 \hat H_0 \epsilon \delta} = \Theta \left( \frac{\gamma}{\delta \epsilon} \right).
		%    \end{align*}
	%    From \textbf{IV}, we set 
	%    \begin{align*}
		%        S = \left( \frac{\beta L}{2 \alpha L_0} \right)^2 = \left( \frac{L}{12 \epsilon L_0}  \right)^2 = \Theta \left( \frac{L^2}{\epsilon^2} \right).
		%    \end{align*}
	%    Note that \textbf{III} and \textbf{IV} together implies Assumption \ref{asm:FO} since
	%    \begin{align*}
		%        \E_{\xi \sim \Xi} \Vert \vg_{T;\mU} (\vx;\xi) \Vert^2 = \Vert \nabla f_{T;\mU}(\vx) \Vert^2 + {\rm Var}_{\xi \sim \Xi} [\vg_{T;\mU} (\vx;\xi)] \le L^2.
		%    \end{align*}
	% Finally, Theorem \ref{thm:hard-T2} indicates that any quantum algorithm requires at least 
	% \begin{align*}
		%     \frac{TS}{2} = \Omega \left( \frac{\gamma}{\delta \epsilon} \left(1 + \frac{L^2}{\epsilon^2} \right)  \right)
		% \end{align*}
	% QSFO complexity to find a $(\delta,\epsilon)$-Goldstein stationary point. Finally, from \textbf{V} we get the condition
	% \begin{align*}
		%     L \ge \left( \frac{\alpha \hat L_0 \sqrt{2T}}{\beta} \right)^2  = \left( 6 \epsilon \hat L_0 \sqrt{2 T} \right)^2 = \frac{2\epsilon \hat L_0 \gamma}{\delta \gamma_0},
		% \end{align*}  
	% which justifies the condition $\delta \ge 16 \epsilon \gamma /L$.
	% \end{proof}

% Indeed, our above proof subsumes the result of \citep[Theorem 3.6]{zhang2023quantum} of the lower bound for finding stationary points in the smooth setting, which the tightest dependency on all parameters.

% \begin{asm}  \label{asm:smooth-SFO}
	% Suppose the objective $f:\sR^d \rightarrow \sR$ satisfies (a) $f$ is $H$-smooth (b) $f^* = \inf_{\vx \in \sR^d} f(\vx) > -\infty$ and $f(\vzero) -  f^* \le \gamma$. (c) we have access to a stochastic gradient estimator $\vg(\vx;\xi)$ such that $\E_{\xi \sim \Xi} [g(\vx;\xi)] = \nabla f(\vx)$ and ${\rm Var}_{\xi \sim \Xi} [\vg(\vx;\xi)] \le \sigma^2$.
	% \end{asm}

% \begin{thm} \label{thm:lb-smooth}
	%  For any $\epsilon>0$, $\gamma > 0 $, $\sigma>\epsilon$, when $d = \tilde \Omega(\gamma H \sigma^2 \epsilon^{-4})$
	%  there exists a 
	% distribution over functions $f: \sR^d \rightarrow \sR$ and a distribution of  
	%     stochastic gradients that satisfy Assumption \ref{asm:smooth-SFO}, such that any quantum algorithm provided with a randomly selected QSFO requires $\Omega(\gamma H \sigma^2 \epsilon^{-4})$ complexity to identify an $(\delta,\epsilon)$-Goldstein stationary point.
	% \end{thm}


% \begin{proof}
	% We use the same hard instance as we prove Theorem \ref{thm:lower-bound} which is $f_{T;\mU} = \alpha \hat f_{T; \mU}(\vx / \beta)$, but use a different scaling on the parameters.  Similar to \textbf{I}, we can set $\alpha = 2 \epsilon \beta$ to give a lower bound for finding $\epsilon$-stationary points. Then, we scale that function to be $H$-smooth by setting $(\alpha \hat H_0 / \beta^2) = H$ as we have done in \textbf{II}. It gives that $\beta = 2 \epsilon \hat H_0 / H $ and thus $\alpha = 4 \epsilon^2 \hat H_0 / H$. Next, we set $T = \gamma / (\gamma_0 \alpha)$ to fulfill to condition of $f_{T; \mU}(\vzero) - \inf_{\vx \in \sR^d} f(\vx) \le \gamma$ as as have done in \textbf{III}. This further yields 
	% \begin{align*}
		%     T = \frac{\gamma}{\gamma_0 \alpha} = \frac{4 H \gamma}{\hat H_0 \epsilon^2} = \Theta \left( \frac{\gamma H}{\epsilon^2} \right).
		% \end{align*}
	% Lastly, we fulfill the bounded variance condition as we have done in \textbf{IV}. We need to ensure $\beta \ge 2 \alpha L_0 \sqrt{S} / \sigma$. It can be satisfied by setting 
	% \begin{align*}
		%     S = \gO \left( \frac{\beta \sigma}{2 \alpha L_0} \right)^2 = \gO \left( \frac{\sigma}{4 \epsilon L_0} \right)^2 = \Theta \left( \frac{\sigma^2}{\epsilon^2} \right). 
		% \end{align*}
	% And the final lower bound of QSFO complexity is
	% \begin{align*}
		%     \frac{TS}{2} = \Omega\left( \frac{\gamma H \sigma^2}{\epsilon^4} \right).
		% \end{align*}
	% \end{proof}
% The above result, strictly match \citep[Theorem 1]{arjevani2023lower} for the classical lower bound, and it is tighter than the quantum lower bound \citep[Theorem 3.6]{zhang2023quantum} of $ \Omega ( (\gamma H \wedge \sigma^2)^2 \epsilon^{-4} )$ in the parameters $\gamma, H$ and $\sigma$.

% We start by showing that $\mA_K$ 
% This function satisfies that
% \begin{enumerate}
	%     \item $\bar f_T(\vzero) = 0$ and $\inf_{\vx \in \sR^T} \bar f_T(\vx) \ge \gamma_0 T $ for $\gamma_0 = 12$. 
	%     \item $\nabla \bar f_T(\vx)$ i $H_0$-Lipschitz with $H_0 = 152$. 
	%     \item For all $\vx \in \sR^T$, $\Vert \nabla \bar f_T(\vx) \Vert_{\infty} \le G_0$ with $G_0 = 23$.
	%     \item For all $\vx \in \sR^T$, 
	%     \begin{align*}
		%         \sqrt{\sum_{i=t}^T x_i^2} \le \frac{1}{2} \quad \Rightarrow \quad \nabla \bar f_T (x_1,x_2,\cdots,x_T) = \nabla \bar f_T (x_1,\cdots,x_t, 0, \cdots, 0).
		%     \end{align*}
	% \end{enumerate}
